% Options for packages loaded elsewhere
\PassOptionsToPackage{unicode}{hyperref}
\PassOptionsToPackage{hyphens}{url}
%
\documentclass[
]{article}
\usepackage{amsmath,amssymb}
\usepackage{iftex}
\ifPDFTeX
  \usepackage[T1]{fontenc}
  \usepackage[utf8]{inputenc}
  \usepackage{textcomp} % provide euro and other symbols
\else % if luatex or xetex
  \usepackage{unicode-math} % this also loads fontspec
  \defaultfontfeatures{Scale=MatchLowercase}
  \defaultfontfeatures[\rmfamily]{Ligatures=TeX,Scale=1}
\fi
\usepackage{lmodern}
\ifPDFTeX\else
  % xetex/luatex font selection
\fi
% Use upquote if available, for straight quotes in verbatim environments
\IfFileExists{upquote.sty}{\usepackage{upquote}}{}
\IfFileExists{microtype.sty}{% use microtype if available
  \usepackage[]{microtype}
  \UseMicrotypeSet[protrusion]{basicmath} % disable protrusion for tt fonts
}{}
\makeatletter
\@ifundefined{KOMAClassName}{% if non-KOMA class
  \IfFileExists{parskip.sty}{%
    \usepackage{parskip}
  }{% else
    \setlength{\parindent}{0pt}
    \setlength{\parskip}{6pt plus 2pt minus 1pt}}
}{% if KOMA class
  \KOMAoptions{parskip=half}}
\makeatother
\usepackage{xcolor}
\usepackage{color}
\usepackage{fancyvrb}
\newcommand{\VerbBar}{|}
\newcommand{\VERB}{\Verb[commandchars=\\\{\}]}
\DefineVerbatimEnvironment{Highlighting}{Verbatim}{commandchars=\\\{\}}
% Add ',fontsize=\small' for more characters per line
\newenvironment{Shaded}{}{}
\newcommand{\AlertTok}[1]{\textcolor[rgb]{1.00,0.00,0.00}{\textbf{#1}}}
\newcommand{\AnnotationTok}[1]{\textcolor[rgb]{0.38,0.63,0.69}{\textbf{\textit{#1}}}}
\newcommand{\AttributeTok}[1]{\textcolor[rgb]{0.49,0.56,0.16}{#1}}
\newcommand{\BaseNTok}[1]{\textcolor[rgb]{0.25,0.63,0.44}{#1}}
\newcommand{\BuiltInTok}[1]{\textcolor[rgb]{0.00,0.50,0.00}{#1}}
\newcommand{\CharTok}[1]{\textcolor[rgb]{0.25,0.44,0.63}{#1}}
\newcommand{\CommentTok}[1]{\textcolor[rgb]{0.38,0.63,0.69}{\textit{#1}}}
\newcommand{\CommentVarTok}[1]{\textcolor[rgb]{0.38,0.63,0.69}{\textbf{\textit{#1}}}}
\newcommand{\ConstantTok}[1]{\textcolor[rgb]{0.53,0.00,0.00}{#1}}
\newcommand{\ControlFlowTok}[1]{\textcolor[rgb]{0.00,0.44,0.13}{\textbf{#1}}}
\newcommand{\DataTypeTok}[1]{\textcolor[rgb]{0.56,0.13,0.00}{#1}}
\newcommand{\DecValTok}[1]{\textcolor[rgb]{0.25,0.63,0.44}{#1}}
\newcommand{\DocumentationTok}[1]{\textcolor[rgb]{0.73,0.13,0.13}{\textit{#1}}}
\newcommand{\ErrorTok}[1]{\textcolor[rgb]{1.00,0.00,0.00}{\textbf{#1}}}
\newcommand{\ExtensionTok}[1]{#1}
\newcommand{\FloatTok}[1]{\textcolor[rgb]{0.25,0.63,0.44}{#1}}
\newcommand{\FunctionTok}[1]{\textcolor[rgb]{0.02,0.16,0.49}{#1}}
\newcommand{\ImportTok}[1]{\textcolor[rgb]{0.00,0.50,0.00}{\textbf{#1}}}
\newcommand{\InformationTok}[1]{\textcolor[rgb]{0.38,0.63,0.69}{\textbf{\textit{#1}}}}
\newcommand{\KeywordTok}[1]{\textcolor[rgb]{0.00,0.44,0.13}{\textbf{#1}}}
\newcommand{\NormalTok}[1]{#1}
\newcommand{\OperatorTok}[1]{\textcolor[rgb]{0.40,0.40,0.40}{#1}}
\newcommand{\OtherTok}[1]{\textcolor[rgb]{0.00,0.44,0.13}{#1}}
\newcommand{\PreprocessorTok}[1]{\textcolor[rgb]{0.74,0.48,0.00}{#1}}
\newcommand{\RegionMarkerTok}[1]{#1}
\newcommand{\SpecialCharTok}[1]{\textcolor[rgb]{0.25,0.44,0.63}{#1}}
\newcommand{\SpecialStringTok}[1]{\textcolor[rgb]{0.73,0.40,0.53}{#1}}
\newcommand{\StringTok}[1]{\textcolor[rgb]{0.25,0.44,0.63}{#1}}
\newcommand{\VariableTok}[1]{\textcolor[rgb]{0.10,0.09,0.49}{#1}}
\newcommand{\VerbatimStringTok}[1]{\textcolor[rgb]{0.25,0.44,0.63}{#1}}
\newcommand{\WarningTok}[1]{\textcolor[rgb]{0.38,0.63,0.69}{\textbf{\textit{#1}}}}
\setlength{\emergencystretch}{3em} % prevent overfull lines
\providecommand{\tightlist}{%
  \setlength{\itemsep}{0pt}\setlength{\parskip}{0pt}}
\setcounter{secnumdepth}{-\maxdimen} % remove section numbering
\ifLuaTeX
  \usepackage{selnolig}  % disable illegal ligatures
\fi
\IfFileExists{bookmark.sty}{\usepackage{bookmark}}{\usepackage{hyperref}}
\IfFileExists{xurl.sty}{\usepackage{xurl}}{} % add URL line breaks if available
\urlstyle{same}
\hypersetup{
  hidelinks,
  pdfcreator={LaTeX via pandoc}}

\author{}
\date{}

\begin{document}

\newcommand{\ba}{\begin{aligned}}
\newcommand{\ea}{\end{aligned}}
\newcommand{\bm}{\begin{bmatrix}}
\newcommand{\em}{\end{bmatrix}}
\newcommand{\bc}{\begin{cases}}
\newcommand{\ec}{\end{cases}}
\newcommand{\empty}{\varnothing}
\newcommand{\se}{\subseteq}
\newcommand{\ve}{\varepsilon}
\newcommand{\s}{\sigma}
\newcommand{\0}{\textbf 0}
\newcommand{\1}{\textbf 1}
\newcommand{\I}{\textbf I}
\newcommand{\X}{\textbf X}
\newcommand{\x}{\textbf x}
\newcommand{\y}{\textbf y}
\newcommand{\b}{\textbf b}
\newcommand{\w}{\textbf w}
\newcommand{\P}{\mathcal P}
\newcommand{\S}{\mathcal S}
\newcommand{\T}{\mathcal T}
\newcommand{\A}{\mathcal A}
\newcommand{\B}{\mathcal B}
\newcommand{\C}{\mathcal C}
\newcommand{\L}{\mathcal L}
\newcommand{\c}{\backslash}
\newcommand{\if}{\text{ if }}
\newcommand{\uinf}[1][k]{\bigcup_{{#1} = 1}^{\infty}}
\newcommand{\sinf}[1][k]{\sum_{{#1} = 1}^{\infty}}
\newcommand{\iinf}[1][k]{\bigcap_{{#1} = 1}^{\infty}}
\newcommand{\plim}{\text{plim }}
\newcommand{\too}{\Longrightarrow}
\newcommand{\sq}{\sqrt}
\newcommand{\dev}[2][i]{({#2}_{#1}-\bar{#2})}
\newcommand{\inv}[1]{{#1}^{-1}}
\newcommand{\add}[1][i]{\sum_{#1 = 1}^n}
\newcommand{\ply}[1][i]{\prod_{#1 = 1}^n}
\newcommand{\br}[1]{\left[#1\right]}
\newcommand{\pr}[1]{\left(#1\right)}
\newcommand{\brace}[1]{\left\{#1\right\}}
\newcommand{\norm}[2][1]{\|{#2}\|_{#1}}
\newcommand{\seq}[1][x]{{#1}_1,{#1}_2,\cdots,{#1}_{n}}
\newcommand{\seqadd}[1][x]{{#1}_1+{#1}_2+\cdots+{#1}_{n}}
\newcommand{\dto}{\overset{\text d}{\to}}
\newcommand{\asim}{\overset{\text a}{\sim}}
\newcommand{\st}{\text{ s.t. }}
\newcommand{\om}[1]{\mu^*\pr{#1}}
\newcommand{\lm}[1]{\mu\pr{#1}}
\newcommand{\d}{\text{ d}}
\newcommand{\dm}{\text{ d}\mu}
\newcommand{\dl}{\text{ d}\lambda}
\newcommand{\abs}[1]{\left|{#1}\right|}

\hypertarget{riemann-integration}{%
\section{Riemann Integration}\label{riemann-integration}}

\hypertarget{introduction-fourier-series}{%
\subsection{Introduction: Fourier
Series}\label{introduction-fourier-series}}

\[F\pr{x} = \frac{a_0}2+\sum_{n = 1}^\infty\br{a_n\cos nx+b_n\sin nx}\]

\[\bc
a_n = \frac{1}{\pi}\int_{-\pi}^{\pi}F\pr{x}\cos nx\text{d}x\\
b_n = \frac{1}{\pi}\int_{-\pi}^{\pi}F\pr{x}\sin nx\text{d}x
\ec\]

In order to validate such results, the operations of \(\int\) and
\(\sum\) for infinite terms should be interchangeable in computation
order.

Justifying Fourier Analysis/Series under Riemann Integral is crazy.
Alternatively, we would define as Lebesgue\textquotesingle s
Integration:

\[\int\sum^{\infty} = \sum^{\infty} \int\]

\hypertarget{riemann-integral}{%
\subsection{Riemann Integral}\label{riemann-integral}}

\textbf{Definition (Partition)} Suppose \(a,b\in \R\) with \(a<b\). A
partition of \([a,b]\) is a finite list of the form
\(x_0,x_1,\cdots,x_n\), where \(a = x_0 < x_1 < \cdots < x_n = b\).

\textbf{Remark:} A partition do not necessarily need to be an even
partition.

\textbf{Definition (Infimum and Supremum of a Function)} If \(f\) is a
real-valued function, \(A\) is a subset of the domain of \(f\), then

\[\inf_Af:=\inf\brace{f\pr x:x\in A}\\
\sup_Af :=\sup \brace{f\pr x:x\in A}\]

\textbf{Definition (Lower and Upper Riemann Sums)} Suppose
\(f:\br{a,b}\to \R\) is a bounded function and \(P\) is a partition
\(x_0,x_1,\cdots,x_n\) of \(\br{a,b}\).

\begin{itemize}
\item
  The lower Riemann sum \(L\pr{f,P,\br{a,b}}\) is defined as

  \[L\pr{f,P,\br{a,b}} := \sum_{j = 1}^n\pr{x_j-x_{j-1}}\cdot \inf_{\br{x_{j-1},x_j}}f\]
\item
  The upper Riemann sum \(U\pr{f,P,\br{a,b}}\) is defined as

  \[U\pr{f,P,\br{a,b}} : = \sum_{j = 1}^n \pr{x_j-x_{j-1}}\cdot \sup_{\br{x_{j-1},x_j}}f\]
\end{itemize}

The next result states that adjoining more points to a partition
increases the lower Riemann sum and decreases the upper Riemann sum.

\emph{\textbf{Proposition (Inequalities with Riemann Sums)}} Suppose
\(f:\br{a,b}\to \R\) is a bounded function and \(P,P'\) are partitions
of \(\br{a,b}\) such that the listing defining \(P\) is a sublist of the
list defining \(P'\). Then

\[L\pr{f,P,\br{a,b}}\le L\pr{f,P',\br{a,b}}\le U\pr{f,P',\br{a,b}}\le U\pr{f,P,\br{a,b}}\]

\begin{center}\rule{0.5\linewidth}{0.5pt}\end{center}

The following result states that if the function is fixed, then each
lower Riemann sum is less than or equal to each upper Riemann sum.

\emph{\textbf{Proposition (Lower Riemann Sums No More Than Upper Riemann
Sums)}} Suppose \(f:\br{a,b}\to \R\) is a bounded function and \(P,P'\)
are partitions of \(\br{a,b}\). Then

\[L\pr{f,P,\br{a,b}}\le U\pr{f,P',\br{a,b}}\]

\begin{center}\rule{0.5\linewidth}{0.5pt}\end{center}

Now that we have been working with lower and upper Riemann sums, here we
define the lower and upper Riemann integrals.

\textbf{Definition (Lower and Upper Riemann Integrals)} Suppose
\(f:\br{a,b}\to \R\) is a bounded function. The \emph{lower Riemann
integral} \(L\pr{f,\br{a,b}}\) and the \emph{upper Riemann integral}
\(U\pr{f,\br{a,b}}\) of \(f\) are defined by

\[\ba
L\pr{f,\br{a,b}} = \sup_P L\pr{f,P,\br{a,b}}\\
U\pr{f,\br{a,b}} = \inf_P U\pr{f,P,\br{a,b}}
\ea\]

In the definition above, we take the supremum (over all partitions) of
the lower Riemann sums because adjoining more points to a partition
increases the lower Riemann sum and should provide a more accurate
estimate of the area under the graph. Similarly, in the definition
above, we take the infimum (over all partitions) of the upper Riemann
sums because adjoining more points to a partition decreases the upper
Riemann sum and should provide a more accurate estimate of the area
under the graph.

Our intuition suggests that for a partition with only a small gap
between consecutive points, the lower Riemann sum should be a bit less
than the area under the graph, and the upper Riemann sum should be a bit
more than the area under the graph. Intuitively, the lower and upper
Riemann Integrals will converge to each other, or converge to a definite
number, as long as the partition goes more and more finely. Both of them
are approximations of the area under \(f\) on \(\br{a,b}\).

\textbf{Remark:} By definition we immediately have
\(L\pr{f,\br{a,b}}\le U\pr{f,\br{a,b}}\). Instead of choosing between
the lower Riemann integral and the upper Riemann integral, the standard
procedure in Riemann integration is to consider only functions for which
those two quantities are equal. When we take equality to this
relationship, that is, \(L\pr{f,\br{a,b}} = U\pr{f,\br{a,b}}\), we call
\(f\) (a bounded function on a closed bounded interval) Riemann
integrable.

\textbf{Definition (Riemann Integrable, Riemann Integral)} A bounded
function on a closed bounded interval is called \emph{Riemann
integrable} if its lower Riemann integral equals its upper Riemann
integral. If \(f:\br{a,b}\to \R\) is Riemann integrable, then the
\emph{Riemann integral} \(\int_a^b f\) is defined by

\[\int_a^b f = L\pr{f,\br{a,b}} = U\pr{f,\br{a,b}}\]

\begin{center}\rule{0.5\linewidth}{0.5pt}\end{center}

\emph{\textbf{Proposition (Continuity and Riemann Integrability)}} Every
continuous real-valued function on each closed bounded interval is
Riemann integrable.

\emph{\textbf{Proposition (Bounds on Riemann Integral)}} Suppose
\(f:\br{a,b}\to \R\) is Riemann integrable. Then

\[\pr{b-a}\inf_{\br{a,b}}f\le \int_a^b f\le \pr{b-a}\sup_{\br{a,b}}f\]

\hypertarget{why-is-riemann-integral-not-so-good}{%
\subsection{Why is Riemann Integral NOT so
good?}\label{why-is-riemann-integral-not-so-good}}

\begin{enumerate}
\def\labelenumi{\arabic{enumi}.}
\item
  \textbf{Riemann integration does NOT handle functions with (too many)
  discontinuities.}

  Consider Dirichlet\textquotesingle s function \(f:\br{0,1}\to \R\):

  \[f\pr x = \bc
  1,\text{ if }x \text{ is rational}\\
  0,\text{ if }x \text{ is irrational}
  \ec\]

  By definition, for \(\br{a,b}\subseteq \br{0,1}\) with \(a<b\), we
  have

  \[\bc
  \inf_{\br{a,b}}f = 0\\
  \sup_{\br{a,b}}f = 1
  \ec\too \bc
  L\pr{f,\br{0,1}} = 0\\
  U\pr{f,\br{0,1}} = 1
  \ec\too L\pr{f,\br{0,1}}\ne U\pr{f,\br{0,1}}\]

  So we conclude that \(f\) is not Riemann integrable.

  Intuitively, since the set of rational numbers is countable and the
  set of irrational numbers is uncountable, \(f\pr{x}\) should be "not
  too" different from the function \(0\). In this way of reasoning,
  \(f\) should, in some sense, have integral \(0\). However, the Riemann
  integral of \(f\) is not defined.
\item
  \textbf{Riemann integration does not work well with unbounded
  functions.}

  Consider the following integration:

  \[\int_{0}^{1}\frac1{\sq x}\text d x\]

  Using the idea of improper integral, this gives an outcome of \(2\).
  However, say \(x_0,x_1,\cdots,x_n\) is a partition of \(\br{0,1}\),
  then \(\sup_{\br{x_0,x_1}}f = \infty\). Thus if we tried to apply the
  definition of the upper Riemann sum to \(f\), we would have
  \(U\pr{f,P,\pr{0,1}} = \infty\) for every partition \(P\) of
  \(\br{0,1}\).
\item
  \textbf{Riemann integration does NOT work well with pointwise limits.}

  Consider the Dirichlet\textquotesingle s function. \(\mathcal Q\) is
  countable. Let \(r_1,r_2,\cdots\) be the sequence that indexes the
  rational numbers. For each positive integer \(k\), we define
  \(f_k:\br{0,1}\to \R\) such that

  \[f_k\pr x = \bc
   1, \text{ if } x\in \brace{r_1,\cdots,r_k}\\
   0,\text{ otherwise}
  \ec\]

  From Real analysis, we see that \(f_k\) is Riemann integrable, and
  \(\int f_k = 0\) (which is a Riemann integral). However, when \(k\)
  goes to \(\infty\), \(f_k\) converges to the
  Dirichlet\textquotesingle s function, which is in contrast not Riemann
  integral.
\end{enumerate}

Because analysis relies heavily upon limits, a good theory of
integration should allow for interchange of limits and integrals, at
least when the functions are appropri- ately bounded. Based on the three
reasons, we hope to develop a new kind of integration that works well
with the issues unsolved by the Riemann integration. We have to develop
a new way to measure subsets of \(\R\).

\hypertarget{measures}{%
\section{Measures}\label{measures}}

\hypertarget{outer-measure}{%
\subsection{Outer Measure}\label{outer-measure}}

Throughout this semester we will work with "infinity" a lot.
It\textquotesingle s more convenient to think of \(\infty\) as a number.
Therefore, to begin with, we base our analysis on the extended real
line:

\[\br{-\infty.\infty} = \R\cup \brace{-\infty,\infty}\]

The underlying intuition is that, \(\infty\) is recognized as a number
instead of a never-reaching limit here.

The computations of \(\infty\) follows that

\[\ba
\infty+\infty & = \infty\\
\infty\cdot \infty & = \infty\\
\ea\]

However, notice that \(\infty-\infty\), \(\infty/\infty\) and
\(0\cdot \infty\) is generally undefined.

\begin{center}\rule{0.5\linewidth}{0.5pt}\end{center}

Our goal is to have a "measure function" of subsets of \(\R\). The
primary question is that, what are the desirable properties?

Suppose \(\mu:\mathcal P\pr \R = 2^{\R}\to \br{-\infty,\infty}\), with
\(\mu\pr{\R} = \infty\) (where \(\P\pr{\R}\) is the set of all subsets
of \(\R\), notated as \(2^{\R}\), or the power set of \(\R\)). We would
wish \(\mu\) to have the following properties:

\begin{enumerate}
\def\labelenumi{\arabic{enumi}.}
\item
  \(\mu\pr{\empty} = 0\), \(\mu\pr{\brace{x}} = 0\).
\item
  \(\mu\pr{\br{a,b}} = b-a\).
\item
  "Linearity"

  \begin{itemize}
  \item
    \(\mu\pr{A} = \mu\pr{A+x}\), where \(x\in \R\). \\
    \textbf{Definition (Translation)} If \(t\in \R\) and
    \(A\subseteq \R\), then the translation \(t+A\) is defined by

    \[t+A:= \brace{t+a:a\in A}\]

    Such property says that measure is translation invariant.
  \item
    \(\mu\pr{\alpha A} = |\alpha|\cdot \mu\pr{A}\).
  \end{itemize}
\item
  \(\mu\pr A\le \mu\pr B\) if \(A\subseteq B\). That is, measure
  preserves order.
\item
  Countable additivity:

  \[\mu\pr{\bigcup_{i = 1}^{\infty}A_i} = \sum_{i = 1}^{\infty}\mu\pr{A_i} \text{ if }A_i \text{'s are pairwise disjoint.}\]

  \textbf{Remarks}: If we alternatively define additivity only in finite
  meanings, such limit operation may not be applicable.
\end{enumerate}

We would depart with "outer measure", which is at first glance intuitive
and then proves to be a useful tool. We would develop the theoretical
works by checking its properties.

\begin{center}\rule{0.5\linewidth}{0.5pt}\end{center}

\textbf{Definition (Length of Open Interval)} The length \(l\pr{I}\) of
an interval \(I\) is defined by

\[l\pr{I} = \bc
b-a, \text{ if }I = \pr{a,b} \text{ for some }a,b\in \R\text{ with }a<b\\
0, \text{ if } I = \empty \\
\infty, \text{ if } I = \pr{-\infty,a} \text{ or }I = \pr{a,\infty} \text{ for some }a \in \R \\
\infty, \text{ if } I = \pr{-\infty,\infty}
\ec\]

Suppose \(A\subseteq \R\). The size of \(A\) should be at most the sum
of the lengths of a sequence of open intervals whose union contains
\(A\). Taking the infimum of all such sums gives a reasonable definition
of the size of \(A\), denoted \(\om{A}\) and called the outer measure of
\(A\).

\textbf{Definition (Outer Measure)} The \emph{outer measure} \(\om{A}\)
of a set \(A\subseteq \R\) is defined by

\[\om{A} = \inf\brace{\sum_{k = 1}^{\infty}l\pr{I_k}:I_1,I_2,\cdots \text{ are open intervals such that }A\subseteq \bigcup_{k = 1}^{\infty}I_k}\]

Intuitively, \(\mu^*\pr A\) corresponds to the length of the smallest
open covering of \(A\).

The definition of outer measure is well-defined, that is, for all
\(A\subset \R\) the outer measure returns a number. One simple argument
is \(A\subset \bigcup_{n = 1}^{\infty}\pr{-n,n}, \forall A\subset \R\),
so there are at least one element in the operation of \(\inf\), which
makes sense.

\textbf{Property (Countable Sets Have Zero Outer Measure)}

\begin{enumerate}
\def\labelenumi{\arabic{enumi}.}
\item
  Finite sets have outer measure \(0\).

  Suppose \(A = \brace{a_1,\cdots,a_n}\) is a finite set of real
  numbers. Suppose \(\ve>0\). Define a sequence \(I_1,I_2,\cdots\) of
  open intervals by

  \[I_k = \bc
  \pr{a_k-\ve,a+\ve} & \text{if }k\le n\\
  \empty  & \text{if } k>n
  \ec\]

  Then \(I_1,I_2,\cdots\) is a sequence of open intervals whose union
  contains \(A\). Clearly \(\sum_{k = 1}^{\infty}l\pr{I_k} = 2\ve n\).
  Hence \(|A|\le 2\ve n\). Because \(\ve\) is an arbitrary positive
  number, this implies that \(|A| = 0\).
\item
  Every countable subset of \(\R\) has outer measure \(0\).

  By definition,
  \(\mu^*\pr A = \inf\brace{\sum l\pr{I_n},\text{ where } A\subset \cup I_n}\).
  Construct
  \(I_n = \pr{a_n - \dfrac{\ve}{2^n},a_n + \dfrac{\ve}{2^n}}\). Clearly,
  \(A\subseteq \bigcup_{n = 1}^{\infty}I_n\). According to the
  definition,

  \[\ba
  \mu^*\pr{A} & \le \sum_{n = 1}^{\infty}l\pr{I_n}\\
  & = \sum_{n = 1}^{\infty}2\cdot \frac{\ve}{2^n} = \ve\sum_{n  =1}^{\infty}\frac{1}{2^{n-1}} = 2\ve\to 0
  \ea\]

  Therefore, for any \(\ve>0\) small, we have \(0\le \mu^*\pr{A}\le 0\);
  immediately we have \(\mu^*\pr{A} = 0\).

  Meanwhile, since finite sets are countable, so it follows that finite
  sets have outer measure \(0\). The first line of the properties is the
  immediate corollary of the second and more general one.
\end{enumerate}

\textbf{Property (Outer Measure Preserves Order)} Suppose \(A\) and
\(B\) are the subsets of \(\R\) with \(A \subseteq B\). Then
\(|A|\le |B|\).

\emph{Proof} Suppose \(I_1,I_2,\cdots\) is a sequence of open intervals
whose union contains \(B\). Then the union of this sequence of open
intervals also contains \(A\). Hence,

\[\om{A}\le \sum_{k = 1}^{\infty}l\pr{I_k}\]

Taking the infimum over all sequences of open intervals whose union
contains \(B\), we have \(\om{A}\le \om{B}\).

\begin{center}\rule{0.5\linewidth}{0.5pt}\end{center}

We expect that the size of a subset of \(\R\) should not change if the
set is shifted to the right or to the left. The next definition allows
us to be more precise.

\textbf{Definition (Translation)} If \(t\in \R\) and \(A\se \R\), then
the \emph{translation} \(t+A\) is defined by

\[t+A = \brace{t+a:a\in A}\]

Obviously, translation does not change the length of an open interval.
The next result states that translation invariance carries over to outer
measure.

\textbf{Property (Outer Measure is Translation Invariant)} Suppose
\(t\in \R\) and \(A\se \R\). Then

\[\om{t+A} = \om{A}\]

\emph{Proof} Suppose \(I_1,I_2,\cdots\) is a sequence of open intervals
whose union contains \(A\). Then \(t+I_1,t+I_2,\cdots\) is a sequence of
open intervals whose union contains \(t+A\). Thus

\[\om{t+A} \le \sum_{k = 1}^{\infty}l\pr{t+I_k} = \sum_{k = 1}^{\infty}l\pr{I_k}\]

Taking the infimum of the last term over all sequences
\(I_1,I_2,\cdots\) of open intervals whose union contains \(A\), we have
\(\om{t+A}\le \om{A}\).

To get the inequality in the other direction, note that
\(A = -t+\pr{t+A}\), so we have
\(\om{A} = \om{-t+\pr{t+A}} \le \om{t+A}\). Hence \(\om{t+A} = \om{A}\).

\textbf{Property (Countable Subadditivity of Outer Measure)} Suppose
\(A_1,A_2,\cdots\) is a sequence of subsets of \(\R\), then

\[\om{\bigcup_{k = 1}^{\infty}A_k}\le \sum_{k = 1}^{\infty}\om{A_k}\]

\textbf{Remarks:} Countable subadditivity implies finite subadditivity,
meaning that

\[\om{A_1\cup\cdots\cup A_n}\le \om{A_1}+\cdots +\om{A_n}\]

for all \(A_1,\cdots ,A_n\in \R\), because we can take \(A_k = \empty\)
for \(k>n\) in the inequality of countable subadditivity.

\begin{center}\rule{0.5\linewidth}{0.5pt}\end{center}

\textbf{Property (Outer Measure of Closed Bounded Interval)} Suppose
\(a,b\in \R\) with \(a<b\), then

\[\om{\br{a,b}} = b-a\]

\emph{Proof} If \(\ve>0\), then
\(\pr{a-\ve,b+\ve},\empty,\empty,\cdots\) is a sequence of open
intervals whose union contains \(\br{a,b}\). Thus
\(|\br{a,b}|\le b-a+2\ve\). Because this inequality holds for all
\(\ve>0\), we conclude that \(|\br{a,b}|\le b-a\).

A proof of the inequality in the opposite direction requires that the
completeness of \(\R\) is used in some form. For example, suppose \(\R\)
was a countable set. Then we would have \(|\br{a,b}| = 0\) since the
outer measure of any countable set is \(0\). Thus something deeper is
going on with the ingredients needed to prove that
\(|\br{a,b}|\ge b-a\). The following definition and Heine-Borel theorem
would be useful:

\textbf{Definition (Open Cover; Finite Subcover)} Suppose \(A\se \R\).

\begin{itemize}
\item
  A collection of \(\C\) of open subsets of \(\R\) is called an
  \emph{open cover} of \(A\) if \(A\) is contained in the union of all
  the sets in \(\C\).
\item
  An open cover of \(\C\) of \(A\) is said to have a \emph{finite
  subcover} if \(A\) is contained in the union of some \emph{finite}
  list of sets in \(\C\).
\end{itemize}

\textbf{Example} The collection \(\brace{\pr{k,k+2}:k\in \Z^+}\) is an
open cover of \([2,+\infty)\) because
\([2,+\infty)\subseteq \bigcup_{k = 1}^{\infty}\pr{k,k+2}\). This open
cover does not have a finite subcover because there do not exist
finitely many sets of the form \(\pr{k,k+2}\) whose union contains
\([2,+\infty)\).

\emph{\textbf{Heine-Borel Theorem}} Every open cover of a \textbf{closed
bounded} subset of \(\R\) has a finite subcover.

\emph{Proof (Continued)} To prove the inequality in the other direction,
suppose \(I_1,I_2,\cdots\) is a sequence of open intervals such that
\(\br{a,b}\se \uinf I_k\). By the Heine-Borel theorem, there exists
\(n\in \Z^+\) such that

\[\br{a,b}\se I_1\cup\cdots \cup I_n\]

We will now prove by induction that the inclusion above implies that

\[\sum_{k = 1}^nl\pr{I_k}\ge b-a\]

This will then imply that
\(\sinf l\pr{I_k}\ge \sum_{k = 1}^nl\pr{I_k}\ge b-a\), completing the
proof that \(|\br{a,b}|\ge b-a\).

To get started with our induction, note that in the case of \(n = 1\)
such relationship stands. Now for the induction step: Suppose \(n>1\)
and \(\br{a,b}\se I_1\cup\cdots \cup I_n\) implies
\(\sum_{k = 1}^nl\pr{I_k}\ge b-a\) for all choices of \(a,b\in\R\) with
\(a<b\). Suppose \(I_1,\cdots,I_n,I_{n+1}\) are open intervals such that

\[\br{a,b}\se I_1\cup\cdots \cup I_n\cup I_{n+1}\]

Thus \(b\) is in at least one of the intervals
\(I_1,\cdots,I_n,I_{n+1}\). By rebelling, we can assume that
\(b\in I_{n+1}\). Suppose \(I_{n+1} = \pr{c,d}\). If \(c<a\), then
\(l\pr{I_{n+1}}\ge b-a\) and there is nothing further to prove; thus we
can assume that \(a<c<b<d\). Hence,
\(\br{a,c}\se I_1\cup\cdots \cup I_n\). By our induction hypothesis, we
have \(\sum_{k = 1}^n l\pr{I_k}\ge c-a\). Thus,

\[\ba
\sum_{k = 1}^{n+1}l\pr{I_k} & = \sum_{k = 1}^nl\pr{I_k}+l\pr{I_{n+1}}\\
& \ge \pr{c-a}+\pr{d-c}\\
& = d-a \\
& \ge b-a
\ea\]

\begin{Shaded}
\begin{Highlighting}[]
\NormalTok{which immediately completes the proof. So we conclude that $\textbackslash{}abs\{\textbackslash{}br\{a,b\}\}\textbackslash{}ge b{-}a$.}
\end{Highlighting}
\end{Shaded}

\textbf{Corollary (Nontrivial Intervals are Uncountable)} Every interval
in \(\R\) that contains at least two distinct elements is uncountable.

\emph{Proof} Suppose \(I\) is an interval that contains \(a,b\in \R\)
with \(a<b\). Then

\[\om{I}\ge \om{\br{a,b}} = b-a> 0\]

\begin{Shaded}
\begin{Highlighting}[]
\NormalTok{where the first inequality above holds because outer measure preserves order, and the equality above comes from the property that the outer measure of each open interval equals its length. Because every countable subset of $\textbackslash{}R$ has outer measure $0$, we can conclude that $I$ is uncountable.}
\end{Highlighting}
\end{Shaded}

\hypertarget{measurable-spaces-and-functions}{%
\subsection{Measurable Spaces and
Functions}\label{measurable-spaces-and-functions}}

\emph{\textbf{Proposition (Nonexistence of Extension of Length to All
Subsets of \(\R\))}} There does not exist a function \(\mu\) with all
the following properties:

\begin{enumerate}
\def\labelenumi{\arabic{enumi}.}
\item
  \(\mu\) is a function from the set of subsets of \(\R\) to
  \(\br{0,\infty}\) (\(\mu:\P\pr{\R} = 2^{\R}\to \br{0,\infty}\)).
\item
  \(\mu\pr{I} = l\pr{I}\) for every open interval \(I\) of \(\R\).
\item
  \(\mu\pr{\bigcup_{k = 1}^{\infty}A_k} = \sum_{k = 1}^{\infty}\mu\pr{A_k}\)
  for every disjoint sequence \(A_1,A_2,\cdots\) of subsets of \(\R\).
\item
  \(\mu\pr{t+A} = \mu\pr A\) for every \(A\subseteq \R\) and every
  \(t\in \R\).
\end{enumerate}

The last result shows that we need to give up one of the desirable
properties in our goal of extending the notion of size from intervals to
more general subsets of \(\R\). We cannot give up (2) because the size
of an interval needs to be its length. We cannot give up (3) because
countable additivity is needed to prove theorems about limits. We cannot
give up (4) because a size that is not translation invariant does not
satisfy our intuitive notion of size as a generalization of length.\\
Thus we are forced to relax the requirement in (1) that the size is
defined for all subsets of \(\R\). Experience shows that to have a
viable theory that allows for taking limits, the collection of subsets
for which the size is defined should be \emph{closed under
complementation and closed under countable unions}. Thus we make the
following definition of \(\s\)-algebra.

\hypertarget{ux3c3--algebra}{%
\subsubsection{\texorpdfstring{\(\s\)-Algebra}{\textbackslash s-Algebra}}\label{ux3c3--algebra}}

\textbf{Definition (\(\s\)-Algebra)} Suppose \(X\) is a set and \(\S\)
is a set of subsets of \(X\). Then \(\S\) is called a
\emph{\(\s\)-algebra} on \(X\) if the following three statements are
satisfied:

\begin{enumerate}
\def\labelenumi{\arabic{enumi}.}
\item
  \(\empty\in \S\). (Notice the notation of \(\in\), an element
  inclusion instead of a set inclusion.)
\item
  If \(E\in \S\), then the complement \(X\backslash E\in \S\).
\item
  If \(E_1,E_2,\cdots\) is a sequence of elements in \(\S\), then
  \(\bigcup_{k = 1}^{\infty}E_k\in \S\).
\end{enumerate}

To verify \(\S\) is indeed a \(\s\)-algebra, it is obvious for the first
two bullet points. For the third one, you need to use the result that
the countable union of countable sets is countable.

\textbf{Examples}:

\begin{enumerate}
\def\labelenumi{\arabic{enumi}.}
\item
  Trivial \(\s\)-algebra: \(X\) is any set, \(S = \brace{\empty,X}\).
\item
  Discrete/Full \(\s\)-algebra: \(X\) is any set, \(S = \P\pr{X} = 2^X\)
  (the power set of \(X\), set of all subsets of \(X\)).
\item
  Suppose \(X\) is a set. Then the set of all subsets \(E\) of \(X\)
  such that \(E\) is countable or \(X\backslash E\) is a \(\s\)-algebra
  on \(X\).
\item
  \(X = \brace{a,b,c,d}\),
  \(S = \brace{\empty,\brace{a,b,c,d},\brace{a,b},\brace{c,d}}\).
\end{enumerate}

\textbf{Property (\(\s\)-Algebras are Closed Under Countable
Intersection)} Suppose \(\S\) is a \(\s\)-algebra on a set \(X\). Then

\begin{itemize}
\item
  \(X\in \S\);
\item
  If \(D,E\in \S\), then \(D\cup E\in \S\) and \(D\cap E\in \S\) and
  \(D\backslash E \in \S\);
\item
  If \(E_1,E_2,\cdots\) is a sequence of elements of \(\S\), then
  \(\bigcap_{k = 1}^{\infty}E_k\in \S\).
\end{itemize}

\textbf{Remarks}: Informally, we think of \(\s\)-algebra as a collection
of sets that are \textbf{closed under intersection, union and
complement.}

\hypertarget{measurable-spaces}{%
\subsubsection{Measurable Spaces}\label{measurable-spaces}}

The word \emph{measurable} is used in the terminology below because in
the next section we introduce a size function, called a \emph{measure},
defined on \emph{measurable sets}.

\textbf{Definition (Measurable Space; Measurable Set)} A
\emph{measurable space} is an ordered pair \(\pr{X,S}\) where \(X\) is a
set and \(\S\) is a \(\s\)-algebra on \(X\). An element of \(\S\) is
called an \emph{\(\S\)-measurable} set, or just a \emph{measurable set}
if \(\S\) is clear from the context.

\textbf{Remarks}: In every discussion you should point out which
\(\s\)-algebra you are currently working on.

\textbf{Example}: If \(X = \R\) and \(\S\) is the set of all subsets of
\(\R\) that are countable or have a countable complement, then the set
of rational numbers is \(\S\)-measurable but the set of positive real
numbers is not \(\S\)-measurable.

\begin{center}\rule{0.5\linewidth}{0.5pt}\end{center}

Now that the full \(\s\)-algebra is not a "good" collection of \(\R\),
what is a "good" \(\s\)-algebra? The next result guarantees that there
is a smallest \(\s\)-algebra on a set \(X\) containing a given set
\(\A\) of subsets of \(X\), which will then pave way for the
introduction of Borel set.

\emph{\textbf{Proposition (Smallest \(\s\)-Algebra Containing a
Collection of Subsets)}} Suppose \(X\) is a set and \(\A\) is a set of
subsets of \(X\). Then the intersection of all \(\s\)-algebras on \(X\)
that contain \(\A\) is a \(\s\)-algebra on \(X\)

\textbf{Examples}

\begin{itemize}
\item
  Suppose \(X\) is a set and \(\A\) is the set of subsets of \(X\) that
  consist of exactly one element:

  \[\A = \brace{\brace{x}:x\in X}\]

  Then the smallest \(\s\)-algebra on \(X\) containing \(\A\) is the set
  of all subsets \(E\) of \(X\) such that \(E\) is countable or
  \(X\backslash E\) is countable.
\item
  Suppose \(\A = \brace{\pr{0,1},\pr{0,\infty}}\). Then the smallest
  \(\s\)-algebra on \(\R\) containing \(\A\) is

  \[\brace{\empty,\pr{0,1},\pr{0,\infty},(-\infty,0]\cup [1,\infty),[1,\infty),\pr{-\infty,1},\R}\]
\end{itemize}

\textbf{Definition (Borel Set)} The smallest \(\s\)-algebra on \(\R\)
containing all \emph{open subsets} of \(\R\) is called the collection of
\emph{Borel subsets} of \(\R\). An element of this \(\s\)-algebra is
called a \emph{Borel set}.

\textbf{Remarks}: We have defined the collection of Borel subsets of
\(\R\) to be the smallest \(\s\)-algebra on \(\R\) containing all the
\emph{open subsets} of \(\R\). We could have defined the collection of
Borel subsets of \(\R\) to be the smallest \(\s\)-algebra on \(\R\)
containing all the \emph{open intervals}, because every open subset of
\(\R\) is the union of a sequence of open intervals.

\textbf{Examples}

\begin{enumerate}
\def\labelenumi{\arabic{enumi}.}
\item
  Every closed subset of \(\R\) is a Borel set because every closed
  subset of \(\R\) is the complement of an open subset of \(\R\).
\item
  Every countable subset of \(\R\) is a Borel set because if
  \(B = \brace{x_1,x_2,\cdots}\). then
  \(B = \bigcup_{k = 1}^{\infty}\brace{x_k}\), which is a Borel set
  because each \(\brace{x_k}\) is a closed set.
\item
  Every half-open interval \([a,b)\) (where \(a,b\in \R\)) is a Borel
  set because \([a,b) = \bigcap_{k = 1}^{\infty}\pr{a-\frac1k,b}\).
\item
  If \(f:\R\to \R\) is a function, then the set of points at which \(f\)
  is continuous is the intersection of a sequence of open sets, and thus
  is a Borel set.
\end{enumerate}

\textbf{Remarks:} We may informally say that any subset of \(\R\) that
you can think of is (highly-probably) a Borel set.

\hypertarget{inverse-images}{%
\subsubsection{Inverse Images}\label{inverse-images}}

\textbf{Definition (Inverse Image)} If \(f:X\to Y\) is a function and
\(A\se Y\), then the inverse image of \(A\) (which is a set), denoted as
\(\inv f\pr{A}\), is defined by

\[\inv f\pr{A} = \brace{x\in X:f\pr{x}\in A}\]

\textbf{Example}: Suppose \(f:\R\to \R\), \(f\pr x = x^2\).

\begin{itemize}
\item
  \(\inv f\pr{\pr{1,4}} = \pr{-2,-1}\cup \pr{1,2}\).
\item
  \(\inv f\pr{\pr{-1,-1}} = \pr{-1,1}\), since
  \(\inv f\pr{\pr{-1,-1}} = \inv f\pr{[0,1)}\).
\end{itemize}

\begin{center}\rule{0.5\linewidth}{0.5pt}\end{center}

Inverse images work very well with set operations.

\emph{\textbf{Proposition (Algebra of Inverse Images)}} Suppose
\(f:X\to Y\) is a function. Then

\begin{enumerate}
\def\labelenumi{\arabic{enumi}.}
\item
  \(\inv f\pr{Y\backslash A} = X\backslash \inv f\pr A\), for
  \(A\subseteq Y\). That is, \(\inv f\pr{A^C} = \pr{\inv f\pr A}^C\).
\item
  \(\inv f\pr{\bigcup_{A\in \A} A} = \bigcup_{A\in \A} \inv f\pr{A}\),
  for every set \(\A\) of subsets of \(Y\).
\item
  \(\inv f\pr{\bigcap_{A\in \A} A} = \bigcap_{A\in \A} \inv f\pr{A}\),
  for every set \(\A\) of subsets of \(Y\).
\end{enumerate}

\textbf{Trick} To prove \(A = B\), it is equivalent to

\[\ba
& A = B\\
\iff & A\subseteq B \text{ and }B\subseteq A\\
\iff & \bc
x\in A\too x\in B\\
x\in B\too x\in A
\ec
\ea\]

\emph{\textbf{Proposition (Inverse Image of a Composition)}} Suppose
\(f:X\to Y\) and \(g:Y\to W\). Then

\[\inv{\pr{g\circ f}}\pr A = \inv f\pr{\inv g\pr A}\]

for every \(A\se W\).

\hypertarget{measurable-functions}{%
\subsubsection{Measurable Functions}\label{measurable-functions}}

The next definition tells us which real-valued functions behave
reasonably with respect to a \(\s\)-algebra on their domain.

\textbf{Definition (Measurable Function)} Suppose \(\pr{X,\S}\) is a
measurable space. A function \(f:X\to \R\) is called \(\S\)-measurable
if \(\inv f\pr B\in \S\) for all Borel set \(B\se \R\).

\textbf{Remarks}

\begin{itemize}
\item
  This definition is a special case of a general class of measurable
  function. In general, suppose \(\pr{X,\S}\) is a measurable space, and
  a function \(f:X\to Y\) is called \(\T\)-measurable if
  \(\inv f\pr B\in \S\) for all \(B\in \T\), where \(\T\) is a
  \(\s\)-algebra of \(Y\).
\item
  (In fact, most of the time we will work with real-valued functions,
  but the works can be extended to more general settings. Real-valued
  functions are the simplest to study with.)
\end{itemize}

\textbf{Example} Consider a set \(X\).

\begin{itemize}
\item
  If \(\S = \brace{\empty, X}\), then the only \(\S\)-measurable
  functions \(f:X\to \R\) are constant functions.
\item
  If \(S = \P\pr X\), then every function is \(\S\)-measurable.
\item
  \(X = \R\). If
  \(\S = \brace{\empty, \pr{-\infty,0},[{0,\infty}),\R}\), then
  \(f:\R\to \R\) is \(\S\)-measurable if and only if \(f\) is constant
  on \(\pr{-\infty,0}\) and \(f\) is constant on \([0,+\infty)\).
\item
  Suppose \(E\) is a subset of a set \(X\). The \emph{characteristic
  function} of \(E\) is the function \(\chi_E:X\to \R\) defined by
  \(\chi_E\pr{x} = \bc
  1 & \if x\in E\\
  0 & \if x\notin E
  \ec\). Suppose \(\pr{X,\S}\) is a measurable space, \(E\se X\) and
  \(B\se \R\). Then

  \[{\chi_E}^{-1}\pr B = \bc
  E & \if 0\notin B \text{ and }1\in B\\
  X\backslash E & \if 0\in B \text{ and }1\notin B\\
  X & \if 0\in B \text{ and }1\in B\\
  \empty & \if 0\notin B \text{ and }1\notin B\\
  \ec\]

  Thus we see that \(\chi_E\) is an \(\S\)-measurable function if and
  only if \(E\in \S\).
\end{itemize}

\hypertarget{weaker-condition-for-measurability}{%
\paragraph{Weaker Condition for
Measurability}\label{weaker-condition-for-measurability}}

The definition of an \(\S\)-measurable function requires the inverse
image of every Borel subset of \(\R\) to be in \(\S\). The following
useful result shows that to verify that a function is \(\S\)-measurable,
we can check the inverse images of a \emph{much smaller} collection of
subsets of \(\R\).

\emph{\textbf{Proposition (Condition for Measurable Function)}} Suppose
\(\pr{X,\S}\) is a measurable space, and \(f:X\to \R\) is a function
that satisfies:

\[\inv f\pr{\pr{a,\infty}} = \brace{x\in X:f\pr{x}>a} \in \S\]

for all \(a\in \R\). Then \(f\) is an \(\S\)-measurable function.

\emph{\textbf{Proof}} Recall the properties of \(\inv f\):

\begin{itemize}
\item
  \(\inv f \pr{\empty} = \empty\).
\item
  If \(A\se \R\), then
  \(\inv f\pr{\R\backslash A} = X\backslash \inv f\pr{A}\).
\item
  If \(A_1,A_2,\cdots\se \R\), then
  \(\inv f \pr{\uinf[i]A_i} = \uinf[i]\inv f\pr{A_i}\).
\end{itemize}

In some sense, the three properties correspond to those of
\(\s\)-algebra. Define

\[\mathcal T = \brace{A\se \R:\inv f\pr A\in \S}\]

We hope to show that every Borel subset of \(\R\) is in \(\mathcal T\).
To achieve this, we will first show that \(\mathcal T\) is a
\(\s\)-algebra on \(\R\).

First because \(\inv f\pr \empty = \empty \in \S\), we have
\(\empty\in \mathcal T\).

If \(A\in \T\), then \(\inv f\pr A\in \S\); hence

\[\inv f\pr{\R\backslash A} = X\backslash \inv f\pr A\in \S\]

The last element inclusion relationship stands since \(\S\) is a
\(\s\)-algebra. Therefore, \(\R\backslash A\in \T\). In other words,
\(\T\) is closed under complementation.

If \(A_1,A_2,\cdots \in \T\), then
\(\inv f\pr{A_1},\inv f\pr{A_2},\cdots \in \S\); hence

\[\inv f\pr{\uinf A_k} = \uinf \inv f\pr{A_k}\in \S\]

The last element inclusion relationship stands since \(\S\) is a
\(\s\)-algebra. Therefore, \(\uinf A_k\in \T\). In other words, \(\T\)
is closed under countable unions.

Based on the above three properties, \(\T\) is a \(\s\)-algebra on
\(\R\).

We note that

\begin{itemize}
\item
  \((a,b] = \pr{a,\infty}\backslash \pr{b,\infty}\), for all \(a<b\).
\item
  \(\pr{a,b} = \uinf (a,b-\frac1k]\).
\end{itemize}

By hypothesis, \(\T\) contains \(\brace{\pr{a,\infty}:a\in \R}\).
Because \(\T\) is closed under complementation, \(\T\) also contains
\(\brace{(-\infty,b]:b\in \R}\). Because the \(\s\)-algebra \(\T\) is
closed under finite intersections, we see that \(\T\) contains
\(\brace{(a,b]:a,b\in \R}\). Because \(\pr{a,b} = \uinf (a,b-\frac1k]\)
and \(\pr{-\infty,b} = \uinf(-k,b-\frac1k]\) and \(\T\) is closed under
countable unions, we can conclude that \(\T\) contains every open subset
of \(\R\). Thus, the \(\s\)-algebra \(\T\) contains the smallest
\(\s\)-algebra on \(\R\) that contains all open subsets of \(\R\). In
other words, \(\T\) contains every Borel subset of \(\R\). Thus \(f\) is
an \(\S\)-measurable function.

By construction we see that, in the result above, we could replace the
collection of sets \(\brace{\pr{a,\infty}:a\in \R}\) by any collection
of subsets of \(\R\) such that the smallest \(\s\)-algebra containing
that collection contains the Borel subsets of \(\R\).

\hypertarget{borel-measurable-functions}{%
\paragraph{Borel Measurable
Functions}\label{borel-measurable-functions}}

We have been dealing with \(\S\)-measurable functions from \(X\) to
\(\R\) in the context of an arbitrary set \(X\) and a \(\s\)-algebra
\(\S\) on \(X\). An important special case of this setup is when \(X\)
is a Borel subset of \(\R\) and \(\S\) is the set of Borel subsets of
\(\R\) that are contained in \(X\). In this special case, the
\(\S\)-measurable functions are called \emph{Borel measurable}.

\textbf{Definition (Borel Measurable Function)} Suppose \(X\se \R\) is a
Borel set. A function \(f:X\to \R\) is called \emph{Borel measurable} if
\(\inv f\pr B\) is a Borel set for every Borel set \(B\se \R\).

\textbf{Remarks}

\begin{itemize}
\item
  If \(X\se \R\) and there exists a Borel measurable function
  \(f:X\to \R\), then \(X\) must be a Borel set since
  \(X = \inv f\pr{\R}\).
\item
  If \(X\se \R\) and \(f:X\to \R\) is a function, then \(f\) is a Borel
  measurable function if and only if \(\inv f\pr{\pr{a,\infty}}\) is a
  Borel set for every \(a\in \R\).
\end{itemize}

\hypertarget{continuity-and-borel-measurability}{%
\paragraph{Continuity and Borel
Measurability}\label{continuity-and-borel-measurability}}

The next result states that continuity interacts well with the notion of
Borel measurability.

\emph{\textbf{Proposition (Every Continuous Function is Borel
Measurable)}} Every continuous real-valued function defined on a Borel
subset of \(\R\) is a Borel measurable function.

\emph{\textbf{Proof}} Suppose \(X\se \R\) is a Borel set and
\(f:X\to \R\) is continuous. To prove that \(f\) is Borel measurable,
fix \(a\in \R\). If \(x\in X\) and \(f\pr x>a\), then by the continuity
of \(f\), there exists some \(\delta_x>0\) such that \(f\pr y>a\) for
all \(y\in \pr{x-\delta_x,x+\delta_x}\cap X\). Thus

\[\inv f \pr{\pr{a,\infty}} = \pr{\bigcup_{x\in \inv f\pr{\pr{a,\infty}}}\pr{x-\delta_x,x+\delta_x}}\cap X\]

The union inside the large parentheses above is an open subset of
\(\R\); hence its intersection with \(X\) is a Borel set. Thus we can
conclude that \(\inv f\pr{\pr{a,\infty}}\) is a Borel set. From the
previous result we conclude that \(f\) is a Borel measurable function.

\textbf{Remarks}: Another way to see this result is using the theorem
that, a function \(f\) is continuous if and only if the inverse image of
an open set is an open set. An immediate corollary is that (combined
with the properties of inverse image), the inverse image of a Borel set
is a Borel set, which completes our proof.

\hypertarget{increasing-function-and-borel-measurability}{%
\paragraph{Increasing Function and Borel
Measurability}\label{increasing-function-and-borel-measurability}}

Next we come to another class of Borel measurable functions. A similar
definition could be made for decreasing functions, with a corresponding
similar result.

\textbf{Definition (Increasing)} Suppose \(X\se \R\) and \(f:X\to \R\)
is a function.

\begin{itemize}
\item
  \(f\) is called \emph{increasing} if \(f\pr x\le f\pr y\) for all
  \(x,y\in X\) with \(x<y\).
\item
  \(f\) is called \emph{strictly increasing} if \(f\pr x<f\pr y\) for
  all \(x,y\in X\) with \(x<y\).
\end{itemize}

\emph{\textbf{Proposition (Every Increasing Function is Borel
Measurable)}} Every increasing function defined on a Borel subset of
\(\R\) is a Borel measurable function.

\emph{\textbf{Proof}} Suppose \(X\se \R\) is a Borel set and
\(f:X\to \R\) is increasing. To prove that \(f\) is Borel measurable,
fix \(a\in \R\). Let \(b = \inf\inv f\pr{\pr{a,\infty}}\), then
\(\inv f\pr{\pr{a,\infty}} = \pr{b,\infty}\cap X\), or
\(\inv f\pr{\pr{a,\infty}} = [b,\infty)\cap X\). Either way, we can
conclude that \(\inv f\pr{\pr{a,\infty}}\) is a Borel set, which implies
that \(f\) is a Borel measurable function.

\textbf{Remarks:} We use the result that, since \(f\) is an increasing
function, \(b = \inf\inv f\pr{\pr{a,\infty}}\) is well-defined; that is,
the infimum always exists.

\hypertarget{composition-and-measurability}{%
\paragraph{Composition and
Measurability}\label{composition-and-measurability}}

The next result shows that measurability interacts well with
composition. (Notice that here we are back discussing the general case
of measurability.)

\emph{\textbf{Proposition (Composition of Measurable Functions is
Measurable)}} Suppose \(\pr{X,\S}\) is a measurable space, and
\(f:X\to \R\) is an \(\S\)-measurable function. Suppose \(g\) is a
real-valued Borel measurable function defined on a subset of \(\R\) that
includes the range of \(f\). Then \(g\circ f:X\to \R\) is
\(\S\)-measurable.

\emph{\textbf{Proof}} Suppose \(B\in \R\) is a Borel set. Then using the
property of inverse image of composition:

\[\pr{g\circ f}^{-1}\pr B = \inv f\pr{\inv g\pr B}\]

Because \(g\) is a Borel measurable function, \(\inv g\pr B\) is a Borel
subset of \(\R\). Because \(f\) is an \(\S\)-measurable function,
\(\inv f\pr{\inv g\pr B}\in \S\). Thus the equation above implies that
\(\pr{g\circ f}^{-1}\pr B\in \S\). Thus \(g\circ f\) is an
\(\S\)-measurable function.

\textbf{Corollary} If \(f\) and \(g\) are continuous, then \(g\circ f\)
is measurable.

\textbf{Example} If \(f\) is measurable, then so are
\(-f,\frac12 f,|f|,f^2,\exp\pr f\), etc. (Because each of these
functions can be written as the composition of \(f\) with a continuous
(and thus Borel measurable) function \(g\).)

\hypertarget{algebraic-operations-and-measurability}{%
\paragraph{Algebraic Operations and
Measurability}\label{algebraic-operations-and-measurability}}

Measurability also interacts well with algebraic operations.

\emph{\textbf{Proposition (Algebraic Operations with Measurable
Functions)}} Suppose \(\pr{X,\S}\) is a measurable space and
\(f,g:X\to \R\) are \(\S\)-measurable. Then

\begin{itemize}
\item
  \(f+g\), \(f-g\), and \(fg\) are \(\S\)-measurable functions;
\item
  If \(g\pr x\ne 0\) for all \(x\in X\), then \(\dfrac{f}{g}\) is an
  \(\S\)-measurable function.
\end{itemize}

\emph{\textbf{Proof}} In order to prove \(f+g\) is an \(\S\)-measurable
function, by definition we are to show for any \(a\in \R\):

\[\pr{f+g}^{-1}\pr{\pr{a,\infty}}\in \S\]

The intuition begins with the observation that, since both \(f\) and
\(g\) are \(\S\)-measurable, if we can find a way (of countable unions)
to represent the inverse image of \(f+g\), then we are done.

\[\pr{f+g}^{-1}\pr{\pr{a,\infty}} = \bigcup_{r\in \Q}\pr{\inv f\pr{\pr{r,\infty}}\cup\inv g\pr{\pr{a-r,\infty}}}\]

In order to prove the equivalence of sets, we use the trick:
\(A= B\iff A\se B \text{ and }B\se A\).

If
\(x\in \bigcup_{r\in \Q}\pr{\inv f\pr{\pr{r,\infty}}\cup\inv g\pr{\pr{a-r,\infty}}}\),
this implies that there exists an \(r\in \Q\) such that
\(x\in\inv f\pr{\pr{r,\infty}}\cup\inv g\pr{\pr{a-r,\infty}}\), since
\(x\) must belong to at least one of the intervals. The immediate result
is that \(x\in \inv f\pr{\pr{r,\infty}}\) and
\(x\in \inv g\pr{\pr{a-r,\infty}}\). By definition of inverse images, we
have \(f\pr{x}\in \pr{r,\infty}\) and \(g\pr x\in \pr{a,r,\infty}\).
Thus, it follows that \(f\pr x+g\pr x\in \pr{a,\infty}\).

Suppose \(x\in \pr{f+g}^{-1}\pr{\pr{a,\infty}}\); that is,
\(f\pr x+g\pr x>a\iff f\pr x > a-g\pr x\). By the completeness of
\(\R\), there must exist some \(r\in \Q\) such that
\(f\pr{x}>r>a-g\pr x\). Therefore, \(x\) belongs to the set in the
right-hand side.

Based on the two directions of intersection relationships we have
proved, the equivalence of the two sets is then proved. From the
intuition we have introduced, the remainder of the proof is trivial.

\begin{itemize}
\item
  By property of composition and measurability, \(-g\) is an
  \(\S\)-measurable function. Thus \(f-g = f+\pr{-g}\) is an
  \(\S\)-measurable function.
\item
  Since \(fg = \dfrac{\pr{f+g}^2-f^2-g^2}{2}\), combined with all the
  properties that we have proved above, the proof is trivial.
\item
  Suppose \(g\pr{x}\ne 0,\forall x\in X\). The function defined on
  \(\R\c \brace{0}\) (which is a Borel subset of \(\R\)) that takes
  \(x\) to \(\frac1x\) is continuous and thus is a Borel measurable
  function. It follows that \(\frac1g\) is an \(\S\)-measurable
  function. Combining this result with what we have already proved about
  the product of \(\S\)-measurable functions, we conclude that
  \(\frac{f}{g}\) is an \(\S\)-measurable function.
\end{itemize}

\hypertarget{limit-and-measurability}{%
\paragraph{Limit and Measurability}\label{limit-and-measurability}}

The next result shows that the pointwise limit of a sequence of
\(\S\)-measurable functions is \(\S\)-measurable. This is a highly
desirable property (recall that the set of Riemann integrable functions
on some interval is not closed under taking pointwise limits).

\emph{\textbf{Proposition (Limit of \(\S\)-Measurable Functions)}}
Suppose \(\pr{X,\S}\) is a measurable space and \(f_1,f_2,\cdots\) is a
sequence of \(\S\)-measurable functions from \(X\) to \(\R\). Suppose
\(\lim_{k\to \infty}f_k\pr x\) exists for each \(x\in X\). Define
\(f:X\to \R\) by

\[f\pr{x}: = \lim_{k\to\infty}f_k\pr{x}\]

Then \(f\) is an \(\S\)-measurable function.

\emph{\textbf{Proof}} In order to prove
\(\inv f\pr{\pr{a,\infty}}\in \S\), we follow the same strategy as
before; that is, to represent \(\inv f\pr{\pr{a,\infty}}\) as a result
of countable unions or intersections.

\[\inv f\pr{\pr{a,\infty}} = \uinf[j]\uinf[m]\bigcap_{k = m}^{\infty}f_k^{-1}\pr{\pr{a+\frac1j,\infty}}\]

Obviously, \(f_k^{-1}\pr{\pr{a+\frac1j,\infty}}\) is an element of
\(\S\). If such equivalence stands, then \(\inv f\pr{\pr{a,\infty}}\) is
also an element of \(\S\), which completes our proof.

\textbf{Put in a Nutshell:}

\begin{itemize}
\item
  Every "well-behaved" function is measurable.
\item
  Algebraic operations preserve measurability.
\item
  \(\lim_{k\to \infty}f_k = f\) is measurable.
\end{itemize}

\hypertarget{working-on-extended-real-line}{%
\subsubsection{Working on Extended Real
Line}\label{working-on-extended-real-line}}

Sometimes it is more convenient for us to work with the \emph{extended
real line}. Hence, we need extension of Borel set for the case of
extended real line.

\textbf{Definition (Borel Subsets of \(\br{-\infty,\infty}\))} A subset
of \(\br{-\infty,\infty}\) is called a \emph{Borel set} if its
intersection with \(\R\) is a Borel set.

\textbf{Remarks:}

\begin{itemize}
\item
  In other words, a set \(C\se \br{-\infty,\infty}\) is a Borel set if
  and only if there exists a Borel set \(B\se \R\) such that \(C = B\)
  or \(C = B\cup\brace{\infty}\) or \(C = B\cup\brace{-\infty}\) or
  \(C = B\cup\brace{-\infty,\infty}\).
\item
  It can be verified that the set of Borel subsets of
  \(\br{-\infty,\infty}\) is a \(\s\)-algebra on
  \(\br{-\infty,\infty}\).
\end{itemize}

\emph{\textbf{Proposition (Weaker Condition for Measurability)}} Suppose
\(\pr{X,\S}\) is a measurable space. A function
\(f:X\to \br{-\infty,\infty}\) is such that
\(\inv f\pr{(a,\infty]}\in \S\) for all \(a \in \R\). Then \(f\) is an
\(\S\)-measurable function.

Such proposition tells us that everything is still good even when we
extend the definition to the extended real line.

\emph{\textbf{Proposition (Infimum and Supremum of a Sequence of
\(\S\)-Measurable Functions are Measurable)}} Suppose \(\pr{X,\S}\) is a
measurable space. \(f_1,f_2,\cdots\) is a sequence of \(\S\)-measurable
functions defined on \(X\to \br{-\infty,\infty}\). Define
\(g,h:X\to\br{-\infty,\infty}\) by

\[\ba
g\pr x &:= \inf\brace{f_k\pr x:k\in \N}\\
h\pr x &:= \sup\brace{f_k\pr x:k\in \N}
\ea\]

Then \(g\) and \(h\) are \(\S\)-measurable functions.

\emph{\textbf{Proof}} Let \(a\in \R\). The definition of supremum
implies that

\[\inv h\pr{(a,\infty]} = \uinf f_k^{-1}\pr{(a,\infty]}\]

The equation above, along with the condition for measurability, implies
that \(h\) is an \(\S\)-measurable function. Then, note that

\[g\pr{x} = -\sup\brace{-f_k\pr{x}:k\in \Z^+},\forall x\in X\]

Thus the result about the supremum implies that \(g\) is an
\(\S\)-measurable function.

\textbf{Remarks:} One of the reason that we hope to work with the
extended real line is that, the \(\inf\) or the \(\sup\) might be
\(\pm \infty\).

\hypertarget{measures-and-their-properties}{%
\subsection{Measures and Their
Properties}\label{measures-and-their-properties}}

\hypertarget{definition}{%
\subsubsection{Definition}\label{definition}}

The original motivation for the next definition came from trying to
extend the notion of the length of an interval. However, the definition
below allows us to discuss size in many more contexts.

\textbf{Definition (Measure)} Let \(\pr{X,\S}\) be a measurable space. A
measure on \(\pr{X,\S}\) is a function \(\mu:\S\to \br{0,\infty}\) such
that

\begin{itemize}
\item
  \(\mu\pr{\empty} = 0\).
\item
  \(\mu\pr{\uinf E_k} = \sinf \mu\pr{E_k}\), for every pairwise disjoint
  sequence \(E_1,E_2,\cdots \in \S\).
\end{itemize}

\textbf{Examples}

\begin{itemize}
\item
  Counting measure: \(X\) is a set, and \(\S = \P\pr X\).

  \[\mu\pr E = \bc
  n \if E \text{ is finite and }E\text{ has }n \text{ elements.}\\
  \infty \if E \text{ is infinite.}
  \ec\]
\item
  \emph{Dirac} measure: \(\pr{X,\S}\) is a measure space. Fix an element
  \(c\in X\).

  \[\delta_c\pr{E} = \bc
  1 \if c\in E\\
  0 \if c\notin E
  \ec\]
\item
  Lebesque measure: \(X = \R, \S = \B\pr \R\). The outer measure is a
  measure.

  Though we have shown that the outer measure is not a measure on the
  power set of \(\R\), since it does not satisfy countable additivity.
  However, it is a measure on \(\pr{\R,\B\pr{\R}}\), as we would show
  later.
\item
  Simple probability measure on a coin toss:
  \(X = \brace{H,T},S = \P\pr{X} = \brace{\empty, \brace{H},\brace{T},\brace{H,T}}\).
  Define \(\mu\) as:

  \[\bc
  \mu\pr{\empty} = 0\\
  \mu\pr{\brace{H}} = \frac12\\
  \mu\pr{\brace{T}} = \frac12\\
  \mu\pr{\brace{H,T}} = 1
  \ec\]

  It can be verified that \(\mu\) is indeed a measure.
\end{itemize}

\begin{center}\rule{0.5\linewidth}{0.5pt}\end{center}

The following terminology is frequently useful.

\textbf{Definition (Measure Space)} \(\pr{X,\S,\mu}\) is called a
measure space. \(X\) is a set, \(\S\) is a \(\s\)-algebra on \(X\), and
\(\mu\) is a measure on the measurable space \(\pr{X,\S}\).

\hypertarget{properties-of-measures}{%
\subsubsection{Properties of Measures}\label{properties-of-measures}}

\emph{\textbf{Proposition (Measure Preserves Order; Measure of a Set
Difference)}} Suppose \(\pr{X,\S,\mu}\) is a measure space. Suppose
\(D,E\in \S\) and \(D\se E\). Then

\begin{itemize}
\item
  \(\mu\pr D\le \mu \pr E\).
\item
  \(\mu\pr{E\c D} = \mu\pr E-\mu\pr D\), given that
  \(\mu\pr{D}<\infty\).
\end{itemize}

\begin{center}\rule{0.5\linewidth}{0.5pt}\end{center}

\emph{\textbf{Proposition (Countable Additivity)}} For any
\(E_1,E_2,\cdots\in \S\), we have

\[\mu\pr{\uinf E_k}\le \sinf \mu\pr{E_k}\]

\textbf{Remarks:} Such countable subadditivity does not require
\(E_1,E_2,\cdots\in \S\) to be pairwise disjoint, which is a general and
useful result.

\emph{\textbf{Proof}} Following the definition of measure, we construct
a sequence of disjoint sets whose union is \(\uinf E_k\). Let
\(D_1 = \empty\) and \(D_k = E_1\cup \cdots \cup E_{k-1}\) for
\(k\ge 2\). Then \(E_1\c D_1,E_2\c D_2,\cdots\) is a disjoint sequence
of subsets of \(X\) whose union equals \(\uinf E_k\). Thus

\[\ba
\lm{\uinf E_k} & = \lm{\uinf\pr{E_k\c D_k}}\\
& = \sinf \lm{E_k\c D_k}\\
& \le \sinf\lm{E_k}
\ea\]

\begin{center}\rule{0.5\linewidth}{0.5pt}\end{center}

Below is the two important properties of the "\emph{continuity of
measure}".

\emph{\textbf{Proposition (Measure of an Increasing Union)}} Suppose
\(\pr{X,\S,\mu}\) is a measure space. Let
\(E_1\se E_2\se E_3\se \cdots\) be an increasing sequence of \(\S\)
measurable sets. Then

\[\mu\pr{\uinf E_k} = \lim_{k\to \infty}\mu\pr{E_k}\]

\emph{\textbf{Proof}} Denote \(E_0 = \empty\) to make notations
consistent. By the construction of an increasing sequence of
\(E_1\se E_2\se E_3\se \cdots\), we can partition \(\uinf E_k\) as:

\[\uinf E_k =\uinf\pr{E_k\c E_{k-1}}\]

Hence,

\[\ba
\mu\pr{\uinf E_k} & = \mu\pr{\uinf \pr{E_k\c E_{k-1}}}\\
& = \sinf \mu\pr{E_k\c E_{k-1}}\\
& = \sinf \pr{\mu\pr{E_k}-\mu\pr{E_{k-1}}}\\
& = \lim_{k \to \infty}\mu\pr{E_k}-\mu\pr{E_0}\\
& = \lim_{k \to \infty}\mu\pr{E_k}
\ea\]

The second equality stands because of countable additivity, and the
third one stands because of the proposition we derived above.

\emph{\textbf{Proposition (Measure of an Decreasing Intersection)}}
Suppose \(\pr{X,\S,\mu}\) is a measure space. Let \(E_1\) be a
decreasing sequence of \(\S\) measurable sets with
\(\mu\pr{E_1}<\infty\). Then

\[\mu\pr{\bigcap_{k = 1}^{\infty}E_k} = \lim_{k\to \infty}\mu\pr{E_k}\]

\emph{\textbf{Proof}} De Morgan\textquotesingle s Law:

\[E_1\c \pr{\bigcap_{k = 1}^{\infty}E_k} = \uinf \pr{E_1\c E_k}\]

Hence,

\[\ba
\mu\pr{E_1\c \pr{\bigcap_{k = 1}^{\infty}E_k}} & = \mu\pr{\uinf \pr{E_1\c E_k}}\\
& = \lim_{k\to\infty}\mu\pr{E_1\c E_k}
\ea\]

Meanwhile, for the left-hand side:

\[\mu\pr{E_1\c \pr{\bigcap_{k = 1}^{\infty}E_k}} = \mu\pr{E_1}-\mu\pr{\bigcap_{k = 1}^{\infty}E_k}\]

For the right-hand side:

\[\lim_{k \to\infty} \mu\pr{E_1\c E_k} = \lim_{k\to \infty}\pr{\mu\pr{E_1}-\mu\pr{E_k}} = \mu\pr{E_1}-\lim_{k\to \infty}\mu\pr{E_k}\]

It is noteworthy that we can subtract \(\mu\pr{E_1}\) from the both
sides of equality only if \(E_1\) is finite. Jointly we have
\(\mu\pr{\bigcap_{k = 1}^{\infty}E_k} = \lim_{k\to \infty}\mu\pr{E_k}\),
which completes the proof.

\textbf{Remarks}: The condition of \(\mu\pr{E_1}<\infty\) can be relaxed
to \(\exist N\in \N: \mu\pr{E_N}<\infty\).

\begin{center}\rule{0.5\linewidth}{0.5pt}\end{center}

\emph{\textbf{Proposition (Measure of a Union)}} Suppose
\(\pr{X,\S,\mu}\) is a measure space. Let \(D,E\in \S\) and
\(\mu\pr{D\cap E}<\infty\). Then

\[\mu\pr{D\cup E} = \mu \pr{D}+\mu\pr{E}-\mu\pr{D\cap E}\]

\emph{\textbf{Proof}} From graphical intuition we separate \(D\cup E\)
into three disjoint parts:

\[D\cup E = \pr{D\c \pr{D\cap E}}\cup \pr{E\c \pr{D\cap E}}\cup \pr{D\cap E}\]

By property of measure, we have

\[\ba
\mu\pr{D\cup E} & = \mu{\pr{D\c \pr{D\cap E}}}+\mu{\pr{E\c \pr{D\cap E}}}+\mu{\pr{D\cap E}}\\
& = \mu\pr{D}-\mu\pr{D\cap E}+\mu\pr{E}-\mu\pr{D\cap E}+\mu\pr{D\cap E}\\
& = \mu \pr{D}+\mu\pr{E}-\mu\pr{D\cap E}
\ea\]

More generally, the \emph{Exclusion - Inclusion Principle} is:

\[\mu\pr{\bigcup_{k = 1}^nE_k} = \sum_{k = 1}^{\infty}\pr{-1}^{k+1}\pr{\sum_{1\le i_1<\cdots < i_k\le n}\mu\pr{E_{i_1}\cap \cdots \cap E_{i_k}}}\]

\hypertarget{lebesgue-measure}{%
\subsection{Lebesgue Measure}\label{lebesgue-measure}}

\hypertarget{additivity-of-outer-measure-on-borel-sets}{%
\subsubsection{Additivity of Outer Measure on Borel
Sets}\label{additivity-of-outer-measure-on-borel-sets}}

The main goal is to show that the "outer measure" \(\mu^*\) is a measure
on \(\pr{\R,\B\pr{\R}}\). (Note that "outer measure" is not itself a
measure; the name is given artificially.) The core task is to answer
whether the "outer measure" satisfies countable additivity. The
procedure unfolds in the following steps:

\begin{itemize}
\item
  \(\mu^*\) is countably additive with open sets.
\item
  \(\mu^*\) is countably additive with closed sets.
\item
  \(\mu^*\) is countably additive with Borel sets.
\end{itemize}

\begin{center}\rule{0.5\linewidth}{0.5pt}\end{center}

\emph{\textbf{Proposition (Additivity of Outer Measure if One of the
Sets is Open)}} Suppose \(A\) and \(G\) are disjoint subsets of \(\R\),
and \(G\) is open. Then

\[\mu^*\pr{A\cup G} = \mu^*\pr A+\mu^* \pr{G}\]

\emph{\textbf{Proof}} The trick we will deploy is the same as before;
that is, to prove \(x = y\) is equivalent to prove both \(x\le y\) and
\(x\ge y\).

Previously, we have proved that the outer measure has countable
subadditivity, so the direction of \(\le\) holds. We will only need to
focus on the opposite direction of \(\ge\). Because of the fact that any
open set can be expressed as a countable union of open intervals, we
first start from the simplest case of \(G\) being an open interval and
then progress to the case where \(G\) is a general open set.

\begin{itemize}
\item
  \(G = \pr{a,b}\) is an open interval with \(a,b\in \R\). We also
  assume that \(a,b\notin A\). Let \(I_1,I_2,\cdots\) be any open
  intervals whose union covers \(A\cup G\). For \(n>0\), let

  \[\ba
  J_n & = I_n\cap \pr{-\infty,a}\\
  K_n & = I_n\cap \pr{a,b}\\
  L_n & = I_n\cap \pr{b,\infty}
  \ea\]

  Note that \(J_n,K_n,L_n\) are themselves open finite intervals. Hence,

  \[l\pr{I_n} = l\pr{J_n} + l\pr{K_n}+l\pr{L_n}\]

  Moreover, we have

  \[\ba
  A& \se \pr{\uinf[n]J_n}\cup \pr{\uinf[n]L_n}\\
  G& \se \uinf[n]K_n
  \ea\]

  By definition of outer measure,

  \[\ba
  & \bc
  \mu^*\pr{A}\le \sinf[n]\pr{l\pr{J_n}+l\pr{L_n}}\\
  \mu^*\pr{G}\le \sinf l\pr{K_n}\\
  \ec\\
  \too & \mu^*\pr{A}+\mu^*\pr{G} \le \sinf[n]\pr{l\pr{J_n}+l\pr{K_n}+ l\pr{L_n}}= \sinf[n]l\pr{I_n}
  \ea\]

  Taking infimum to both sides (over all open coverings of \(A\cup G\)),
  we have

  \[\mu^*\pr{A}+\mu^*\pr{G}\le \mu^*\pr{A\cup G}\]

  \textbf{Remarks:} Here \(G\) is assumed to be open to ensure that
  \(K_n\)\textquotesingle s are open intervals, otherwise some of the
  properties of outer measure are not applicable.
\item
  \(G\) is any open set. (In fact, any open set can always be decomposed
  into the union of disjoint intervals. Equivalently,
  \(G = \uinf[n] I_n\).) Since \(\pr{\bigcup_{n = 1}^N I_n}\se G\), we
  have

  \[\ba
  \mu^*\pr{A\cup G}& \ge \mu^*\pr{A\cup \pr{\bigcup_{n = 1}^N I_n}}\\
  & = \mu^*\pr{A} + \sum_{n = 1}^N\mu^*\pr{I_n}\\
  & = \mu^*\pr{A} + \sum_{n = 1}^Nl\pr{I_n}
  \ea\]

  Note that the second equality is true because of result in the first
  case we have just discussed.

  Taking limit of \(N\to \infty\) of both sides, we have

  \[\ba
  \mu^*\pr{A\cup G}&\ge \mu^*\pr{A}+\sinf[n]l\pr{I_n}\\
  & \ge \mu^*\pr{A}+\mu^*\pr{G}
  \ea\]
\end{itemize}

\begin{center}\rule{0.5\linewidth}{0.5pt}\end{center}

\emph{\textbf{Proposition (Additivity of Outer Measure if One of the
Sets is Closed)}} Suppose \(A\) and \(F\) are disjoint subsets of
\(\R\), and \(F\) is closed. Then

\[\mu^*\pr{A\cup F} = \mu^*\pr{A}+\mu^*\pr{F}\]

\textbf{Remarks:} For \(F\) to be closed, the complement of \(F\) is
open.

\emph{\textbf{Proof}} Suppose \(I_1,I_2,\cdots\) are open intervals
whose union covers \(A\cup F\). Let \(G = \uinf I_k\). Since
\(A\cap F = \empty\), we have \(A\se \pr{G\c F}\), which implies
\(\mu^*\pr{A}\le \mu^*\pr{G\c F}\). Note that \(G\c F = G\cap F^c\), and
both \(G\) and \(F^c\) are open, so \(G\c F\) is open.

\[\ba
 \mu^*\pr{A}& \le \mu^*\pr{G\c F}\\
\too \mu^*\pr{A}+\mu^*\pr{F}&\le \mu^*\pr{G\c F}+\mu^*\pr{F}\\
& = \mu^*\pr{\pr{G\c F}\cup F}\\
& = \mu^*\pr{G}\\
& \le \sinf l\pr{I_k}
\ea\]

Note that the second equality holds because of the proposition we have
just proved above.

Taking infimum of both sides, we get the final result that

\[\mu^*\pr{A}+\mu^*\pr{F}\le \mu^*\pr{A\cup F}\]

\begin{center}\rule{0.5\linewidth}{0.5pt}\end{center}

\emph{\textbf{Proposition (Approximation of Borel Sets by Some Closed
Set)}} Suppose \(B\se \R\) is a Borel set. Then for every \(\ve>0\),
there exists a closed set \(F\se B\) such that

\[\mu^*\pr{B\c F}\le \ve\]

\emph{\textbf{Proof}} To prove something is true for Borel set, you need
to construct a collection of sets that satisfies such property, and show
such collection is

\begin{itemize}
\item
  Closed under intersection/union.
\item
  Closed under complement.
\item
  Contains open sets.
\end{itemize}

Let

\[\mathcal L = \brace{D\se \R |\text{ For }\ve >0, \text{ there exists a closed set }F\se D \text{ such that }\mu^*\pr{D\c F}\le \ve}\]

We hope to show:

\begin{itemize}
\item
  \(\mathcal L\) is itself a \(\s\)-algebra.
\item
  \(\mathcal L\) contains open sets/open intervals.
\end{itemize}

Then based on the two points, \(\mathcal L\) is a \(\s\)-algebra that
contains the Borel \(\s\)-algebra.

It is in fact trivial to see the second point. Let \(D = \pr{a,b}\) with
\(a<b\) and \(a,b\in \R\). Take
\(F = \br{a+\frac{\ve}2,b-\frac{\ve}2}\), and then we are done. The main
contents leave for proving \(\mathcal L\) being a \(\s\)-algebra.

\begin{itemize}
\item
  \(\mathcal L\) contains \(\empty\): If \(D = \empty\), then take
  \(F = \empty\), which is open and then valid.
\item
  \(\mathcal L\) is closed under countable intersection and union.
  However, note that it is equivalent for being closed under countable
  intersection and being closed under countable union.

  Suppose \(D_1,D_2,\cdots \in \L\). Fix \(\ve>0\). For each \(k>0\),
  there exists a closed set \(F_k\se D_k\) such that

  \[\mu^*\pr{D_k\c F_k}\le \frac{\ve}{2^k}\]

  A fact from set theory, that the intersection of closed sets is
  closed, regardless of whether or not such operation is finite or
  infinite. Then \(\iinf F_k\) is a closed set. By definition, since
  each \(F_k\) is a subset of \(D_k\), we have

  \[\ba
  \iinf F_k\se \iinf D_k\\
  \ea\]

  Moreover, you can check it for yourself that

  \[\pr{\iinf D_k}\c \pr{\iinf F_k}\se \uinf \pr{D_k\c F_k}\]

  Then we have

  \[\mu^*\pr{\pr{\iinf D_k}\c \pr{\iinf F_k}}\le \sinf \mu^*\pr{D_k\c F_k}\le \sinf \frac{\ve}{2^k} = \ve\]

  The result shows that \(\iinf D_k\in \L\).
\item
  \(\L\) is closed under complements.

  Suppose \(D\in \L\). We had to show \(D^c\in \L\).

  \begin{itemize}
  \item
    Suppose \(\mu^*\pr{D}<\infty\). Fix \(\ve>0\). Let \(F\se D\) such
    that \(F\) is closed and \(\mu^*\pr{D\c F}\le\frac{\ve}2\). From
    definition of \(\mu^*\pr{D}\), there is an open set \(G\) (which can
    be imagined as the union of an open covering) with \(G\se D\) such
    that

    \[\mu^*\pr{G}\le \mu^*\pr{D}+\frac{\ve}2\]

    First because \(G\) is open, \(G^c\) is closed. Because \(G\se D\),
    \(G^c\se D^c\). And

    \[D^c\c G^c = G\c D\se G\c F\]

    Then

    \[\ba
    \mu^*\pr{D^c\c G^c} & = \mu^*\pr{G\c D} \\
    & \le \mu^*\pr{G\c F} \\
    & = \mu^*\pr{G}-\mu^*\pr{F} \\
    & =\pr{\mu^*\pr{G}-\mu^*\pr{D}}+\pr{\mu^*\pr{D}-\mu^*\pr{F}}\\
    & \le \frac{\ve}2+\mu^*\pr{D\c F}\\
    & \le \frac{\ve}2+\frac{\ve}2 = \ve
    \ea\]

    Note that the equality in the second line comes from additivity of
    outer measure when one of the sets is closed, and the fourth line
    above uses subadditivity applied to the union
    \(D = \pr{D\c F}\cup F\).
  \item
    When \(\om{D} = \infty\), for any \(D\in \L\), consider

    \[D_k = D\cap \br{-k,k},\forall k\in \N\]

    Because \(\L\) is closed under intersection, and both \(D\) and
    \(\br{-k,k}\) are in \(\L\), we have \(D_k\in \L\). From
    construction, we see that \(D_k\) has a finite measure, so
    \(D_k^c\in \L\). It follows that

    \[\ba
    D & = \uinf D_k\\
    \too D^c & = \iinf D_k^c
    \ea\]

    Again, since each \(D_k^c\in \L\), and \(\L\) is closed under
    intersection, so \(D^c\in \L\), which completes our proof.
  \end{itemize}
\end{itemize}

\emph{\textbf{Proposition (Additivity of Outer Measure if One of the
Sets is a Borel Set)}} Suppose \(A\) and \(B\) are two disjoint sets,
and \(B\) is a Borel set. Then

\[\mu^*\pr{A\cup B} = \mu^*\pr{A}+\mu^*\pr{B}\]

\emph{\textbf{Proof}} From the previous proposition, since \(B\) is a
Borel set, there exists a closed set \(F\se B\) such that
\(\mu^*\pr{B\c F}\le\ve\). Since \(F\) is a closed set, again from the
previous results,

\[\ba
\mu^*\pr{A\cup B}& \ge \mu^*\pr{A\cup F} \\
& = \mu^*\pr{A}+\mu^*\pr{F}\\
& \ge \mu^*\pr{A}+\mu^*\pr{B}-\mu^*\pr{B\c F}\\
& \ge \mu^*\pr{A}+\mu^*\pr{B}-\ve
\ea\]

The first inequality holds because of the fact that
\(F\se B\too \pr{A\cup F}\se\pr{A\cup B}\). The inequality in the third
line is true by countable subadditivity of the outer measure. The result
above holds for any \(\ve>0\). Take limit \(\ve\to 0\), we end up with

\[\mu^*\pr{A\cup B}\ge \mu^*\pr{A}+\mu^*\pr{B}\]

\textbf{Remarks:} We should never write
\(\mu^*\pr{B\c F} = \mu^*\pr{B}-\mu^*\pr{F}\) for \(F\se B\), because we
are not sure whether \(B\) and \(F\) are infinite or not. What we always
have is countable subadditivity of the outer measure,

\[\ba
&\ba
B = F\cup\pr{B\c F}& \too \mu^*\pr{B}\le \mu^*\pr{F}+\mu^*\pr{B\c F}
\\F = B\cup \pr{F\c B}&\too \mu^*\pr{F}\le \mu^*\pr{B}+\mu^*\pr{F\c B}
\ea\\
\too &\bc
\mu^*\pr{B\c F}\ge \mu^*\pr{B}-\mu^*\pr{F}\\
\mu^*\pr{F\c B}\ge \mu^*\pr{F}-\mu^*\pr{B}
\ec
\ea\]

\textbf{Corollary (Existence of a Subset of \(\R\) that is Not a Borel
Set)} There exists a set \(B\se \R\) such that \(\mu^*\pr{B}<\infty\)
and \(B\) is not a Borel set.

\emph{Proof} In the section of countable subadditivity of the outer
measure we showed that there exist disjoint sets \(A,B\se \R\) such that
\(\mu^*\pr{A\cup B}\ne \mu^*\pr{A}+\mu^*\pr{B}\). For any such sets, we
must have \(\mu^*\pr{B}<\infty\) because otherwise both
\(\mu^*\pr{A\cup B}\) and \(\mu^*\pr{A}+\mu^*\pr B\) equal \(\infty\)
(as follows from the inequality \(\mu^*\pr{B}\le \mu^*\pr{A\cup B}\)).
Now the additivity of outer measure if one of the sets is a Borel sets
implies that \(B\) is not a Borel set.

\begin{center}\rule{0.5\linewidth}{0.5pt}\end{center}

The tools we have constructed now allow us to prove that outer measure,
when restricted to the Borel sets, is a measure.

\emph{\textbf{Proposition (Outer Measure is a Measure on Borel Sets)}}
The outer measure \(\mu^*\) is a measure on \(\pr{\R,\B\pr{\R}}\).

\emph{\textbf{Proof}} Consider \(A_1,A_2,\cdots\in \B\pr{R}\) are
disjoint. To prove the countable additivity, because we have had
countable subadditivity of the outer measure, we only need to prove

\[\mu^*\pr{\uinf[n]A_n}\ge \sinf[n]\mu^*\pr{A_n}\]

To see this, we first focus on finite cases to employ the preceding
result:

\[\ba
\mu^*\pr{\uinf[n]A_n} & \ge \mu^*\pr{\bigcup_{n = 1}^N A_n}\\
& = \sum_{n = 1}^N \mu^*\pr{A_n}\overset{N\to\infty}{\longrightarrow}\sinf \mu^*\pr{A_n}
\ea\]

The second equality holds because each of \(A_n\) is a Borel set and we
can then apply the previous proposition. Note that the limit here is
valid itself just from our definition of infinite series.

\textbf{Definition (Lebesgue Measure)} "The" Lebesgue measure on
\(\pr{\R,\B\pr{\R}}\) is \(\mu:\B\pr{\R}\to \br{0,\infty}\) such that

\[\mu\pr{B}:=\mu^*\pr{B},\forall B\in \B\pr{\R}\]

\textbf{Remarks:} The actual Lebesgue measure is defined on a "slightly"
bigger \(\s\)-algebra. There are some Lebesgue measurable sets that are
not Borel sets; but that definition is enough to serve our purpose.

\hypertarget{lebesgue-measurable-sets}{%
\subsection{Lebesgue Measurable Sets}\label{lebesgue-measurable-sets}}

As we will see in this subsection, outer measure is actually a measure
on a somewhat larger class of sets called the Lebesgue measurable sets.
The approach chosen here has the advantage of emphasizing that a
Lebesgue measurable set differs from a Borel set by a set with outer
measure \(0\). The attitude here is that sets with outer measure \(0\)
should be considered small sets that do not matter much.

\textbf{Definition (Lebesgue Measurable Set)} A set \(A\se \R\) is
called \emph{Lebesgue measurable} if there exists a Borel set \(B\se A\)
such that the outer measure

\[\mu^*\pr{A\c B} = 0\]

Intuitively, Lebesgue just added some sets that differ from Borel sets
by a set with outer measure of \(0\) into the Borel sets and call them
together the \emph{Lebesgue measurable} sets. Each Borel set is Lebesgue
measurable because if \(A\se \R\) is a Borel set, then we can take
\(B = A\) in the definition above.

\emph{\textbf{Proposition (Different Definitions for Lebesgue
Measurability)}} Suppose \(A\se \R\). Then all of the following
statements are equivalent:

\begin{itemize}
\item
  \(A\) is Lebesgue measurable.
\item
  \textbf{There exists a Borel set \(B\se A\) such that
  \(\mu^*\pr{A\c B} = 0\).}
\item
  For each \(\ve>0\), there exists a \textbf{closed} set \(F\se A\) such
  that

  \[\mu^*\pr{A\c F}<\ve\]
\item
  There exists a (countable) sequence of closed sets \(F_1,F_2,\cdots\)
  such that

  \[\mu^*\pr{A\c \pr{\uinf F_k}} = 0\]
\item
  For each \(\ve>0\) there exists an \textbf{open} set \(G \supseteq A\)
  such that

  \[\mu^*\pr{G\c A}<\ve\]
\item
  There exists a (countable) sequence of closed sets \(G_1,G_2,\cdots\)
  such that

  \[\mu^*\pr{\pr{\iinf G_k}\c A} = 0\]
\end{itemize}

\textbf{Remarks:}

\begin{itemize}
\item
  The fourth condition implies that every Borel set is almost a
  countable union of closed sets. The last condition implies that every
  Borel set is almost a countable intersection of open sets.
\item
  The reason why we have developed such list of definitions is that,
  they would allow us to work on some more general settings (such as
  beyond the \(\R\); e.g., the set of functions, vectors, etc.)
\item
  The infinite union of closed sets is not necessarily a closed set.
  That is why we developed the second and the third definitions (instead
  of regard them as being equivalent). For example,
  \(\uinf[n]\br{\frac1n,1-\frac1n} = \pr{0,1}\).
\item
  Carethedory\textquotesingle s criterion: \(A\) is Lebesgue measurable
  if and only if, for all \(E\se \R\), we have

  \[\mu^*\pr{E} = \mu^*\pr{E\cap A}+\mu^*\pr{E^c\cap A}\]

  The intuition is that, if \(A\) is Lebesgue measurable, it should
  separate each set \(E\in \R\) well. Compared with the previous
  definitions, this definition is even more general; that is, it imposes
  on restriction on \(A\)\textquotesingle s being closed or open.
\end{itemize}

\emph{\textbf{Proposition (Outer Measure is a Measure on Lebesgue
Measurable Sets)}} The collection \(\L\) of Lebesgue measurable sets is
a \(\s\)-algebra on \(\R\). The outer measure is a measure on
\(\pr{\R,\L\pr{\R}}\), which is called the Lebesgue measure.

\emph{\textbf{Proof}} \(\L\) indeed contains both the Borel sets and
every set with outer measure \(0\), which is a \(\s\)-algebra on \(\R\).
To prove the countable additivity, suppose \(A_1,A_2,\cdots\) is a
disjoint sequence of Lebesgue measurable sets. By the definition of
Lebesgue measurable set, for each \(k\in \Z^+\) there exists a Borel set
\(B_k\se A_k\) such that \(\mu^*\pr{A_k\c B_k} = 0\). Now

\[\ba
\mu^*\pr{\uinf A_k}\ge \mu^*\pr{\uinf B_k} = \sinf \mu^*\pr{B_k} = \sinf\mu^*\pr{A_k}
\ea\]

Note that the second line above holds because \(B_1,B_2,\cdots\) is a
disjoint sequence of Borel sets and outer measure is a measure on Borel
sets; the last line above holds because

\[\ba
B_k\se A_k& \too \mu^*\pr{B_k}\le \mu^*\pr{A_k}\\
A_k = B_k\cup \pr{A_k\c B_k}& \too \mu^*\pr{A_k}\le \mu^*\pr{B_k}+\mu^*\pr{A_k\c B_k} = \mu^*\pr{B_k}
\ea\too \mu^*\pr{A_k} = \mu^*\pr{B_k}\]

The inequality above, combined with countable subadditivity of outer
measure, implies that \(\mu^*\pr{\uinf A_k} = \sinf\mu^*\pr{A_k}\),
completing the proof of the outer measure being a measure on
\(\pr{\R,\B\pr{\R}}\).

From now on, when we say "measure" we mean \(\mu^*:\B\pr{\R}\to \R\),
and elements in \(\B\pr{\R}\) are called "measurable sets".

\hypertarget{convergence-of-measurable-functions}{%
\subsection{Convergence of Measurable
Functions}\label{convergence-of-measurable-functions}}

Littlewood\textquotesingle s 3 principles of measure theory:

\begin{enumerate}
\def\labelenumi{\arabic{enumi}.}
\item
  Every measurable set is "nearly" a finite union of intervals.
\item
  (Luzin\textquotesingle s Theorem) Every measurable function is
  "nearly" continuous.
\item
  (Egorov\textquotesingle s Theorem) Every convergent sequence of
  functions is "nearly" uniformly convergent.
\end{enumerate}

\begin{center}\rule{0.5\linewidth}{0.5pt}\end{center}

\textbf{Definition (Pointwise Convergence and Uniform Convergence)}
Suppose \(X\) is a set, \(f_1,f_2,\cdots: X\to \R\) is a sequence of
functions and \(f:X\to \R\) is a function.

\begin{itemize}
\item
  Pointwise convergence: \(f_1,f_2,\cdots\) converges pointwise to \(f\)
  on \(X\) if for both \(x\in X\) and \(\ve>0\), there exists
  \(n\in \N\) such that \(\abs{f_k\pr{x}-f\pr{x}}<\ve\) for all
  \(k\ge n\). (\(\lim_{k\to\infty} f_k\pr{x} - f\pr{x},\forall x\in X\))
\item
  Uniform convergence: \(f_1,f_2,\cdots\) converges uniformly to \(f\)
  on \(X\) if for every \(\ve>0\), there exists \(n\in \N\) such that
  \(\abs{f_k\pr{x}-f\pr{x}}<\ve\) for all \(k\ge n\) and for all
  \(x\in X\).
\end{itemize}

\emph{\textbf{Theorem (Uniform Limit of Continuous Functions is
Continuous)}} Suppose \(B\se \R\) and \(f_1,f_2,\cdots\) is a sequence
of functions from \(B\) to \(\R\) that converges uniformly on \(B\) to a
function \(f:B\to \R\). Suppose \(b\in B\) and \(f_k\) is continuous at
\(b\) for each \(k\in \Z^+\). Then \(f\) is continuous at \(b\).

\hypertarget{egorovs-theorem}{%
\subsubsection{Egorov\textquotesingle s Theorem}\label{egorovs-theorem}}

A sequence of functions that converges pointwise need not converge
uniformly. However, the next result says that a pointwise convergent
sequence of functions on a measure space with finite total measure
almost converges uniformly, in the sense that it converges uniformly
except on a set that can have arbitrarily small measure.

\emph{\textbf{Egorov\textquotesingle s Theorem}} Suppose
\(\pr{X,\S,\mu}\) is a measure space with \(\mu\pr{X}<\infty\). Suppose
\(f_1,f_2,\cdots\) is a sequence of \(\S\)-measurable function from
\(X\) to \(\R\) that converges pointwise to a function \(f:X\to \R\).
Then for every \(\ve>0\), there exists an \(\S\)-measurable set \(E\)
such that \(\mu\pr{X\c E}<\ve\) and \(f_1,f_2,\cdots\) converges
uniformly to \(f\) on \(E\).

\emph{\textbf{Proof}} Let \(\ve>0\). Fix \(n\in\N\). The fact that
\(f_n\to f\) pointwise on \(X\) implies that

\[X = \uinf[m]\bigcap_{k = m}^{\infty}\brace{x\in X: \abs{f_K\pr{x}-f\pr x}<\frac1n}\]

For \(m\in\N\), let

\[A_{m,n} = \bigcap_{k = m}^{\infty}\brace{x\in X: \abs{f_K\pr{x}-f\pr x}<\frac1n}\]

Note that each \(A_{m,n}\in \S\) because each \(f_k-f\) is an
\(\S\)-measurable function. Then
\(A_{1.n}\se A_{2,n}\se A_{3,n}\se \cdots\) is an increasing sequence of
sets. Hence, by continuity of measure, we have

\[\lim_{m\to\infty}\mu\pr{A_{m,n}} = \lm{\uinf[m]A_{m,n}} = \lm{X}\]

The equality means there exists \(m_n\in \N\) (an integer dependent of
\(n\)) such that

\[\lm{X}-\lm{A_{m_n,n}}<\frac{\ve}{2^n}\]

Let

\[E = \iinf[n]A_{m_n,n}\]

Then

\[\ba
\lm{X\c E} & = \lm{X\c \pr{\iinf[n]A_{m_n,n}}}\\
& = \lm{\uinf[n]\pr{X\c A_{m_n,n}}}\\
& \le \sinf[n]\lm{X\c A_{m_n,n}} \\
& \le \sinf[n]\frac{\ve}{2^n} = \ve
\ea\]

where the equality in the second line is by De Morgan\textquotesingle s
law, and the inequality in the third line is by subadditivity of
\(\mu\). (Note that countable additivity is only applicable to the case
of disjoint sets, while subadditivity is true for all interested sets.)

To complete the proof, we must verify that \(f_1,f_2,\cdots\) converges
uniformly to \(f\) on \(E\). To do this, suppose \(\ve'>0\). Let
\(n\in \Z^+\) be such that \(\frac1n<\ve'\). Then \(E\se A_{m_n,n}\),
which implies that

\[\abs{f_k\pr x-f\pr x}<\frac1n<\ve'\]

for all \(k\ge m_n\) and all \(x\in E\). Hence \(f_1,f_2,\cdots\) does
indeed converge uniformly to \(f\) on \(E\).

\textbf{Remarks:}

\begin{itemize}
\item
  Since \(f_1,f_2,\cdots\) are measurable functions, it follows that
  \(A_{m_n,n}\)\textquotesingle s and \(E\) are measurable.
\item
  To prove a set is measurable, try to describe it as a union or
  intersection of sets.
\end{itemize}

\hypertarget{luzins-theorem}{%
\subsubsection{Luzin\textquotesingle s Theorem}\label{luzins-theorem}}

Here we are going to introduce the concept of simple functions, and show
that any function can be approximated by a sequence of simple functions
that is pointwise and uniformly convergent. The results of simple
function would be quite helpful for further works.

\textbf{Definition (Characteristic/Indicator Function)} Let \(E\se X\)
be any subset of \(X\). The characteristic function (or, indicator
function) on \(E\), \(\chi_{E}:X\to \R\) is defined as

\[\chi_E\pr{x} =\bc
1 \if x\in E\\
0\if x \in E
\ec\]

\textbf{Definition (Simple Function)} A function is called \emph{simple}
if it takes on only finitely many values. Let \(\pr{X,\S}\) be a
measurable space. An \(\S\)-measurable simple function \(f:X\to \R\) has
the form

\[f\pr{x} = c_1\chi_{E_1}+c_2\chi_{E_2}+\cdots+c_n\chi_{E_n}\]

for \emph{distinct} nonzero \(c_1,c_2,\cdots,c_n\) and
\(E_1,E_2,\cdots ,E_n\in \S\).

\textbf{Remarks:} Step functions are in analog to simple functions.
However, different from step functions that are defined on intervals,
simple functions are defined on \emph{measurable sets}.

\emph{\textbf{Theorem (Approximation by Simple Functions)}} Suppose
\(\pr{X,\S}\) is a measurable space and \(f:\br{-\infty,\infty}\) is
\(\S\)-measurable. Then there exists a sequence \(f_1,f_2,\cdots\) of
functions from \(X\) to \(\R\) such that

\begin{itemize}
\item
  each \(f_k\) is a simple \(\S\)-measurable function;
\item
  \(\abs{f_k\pr x}\le \abs{f_{k+1}\pr x}\le \abs{f\pr x}\) for all
  \(k\in \N\) and all \(x\in X\);
\item
  \(\lim_{k\to\infty} f_k\pr x = f\pr x\) for every \(x\in X\);
\item
  \(f_1,f_2,\cdots\) converges \emph{uniformly} on \(X\) to \(f\) if
  \(f\) is bounded.
\end{itemize}

\emph{\textbf{Proof}} Instead of directly defining \(f\) on its
intervals (like the Riemann integrals do), we define \(f\) on the
inverse image of evenly-cut intervals of the codomain of \(f\):

For \(k>0\) and \(n\in \Z\), we divide \([n,n+1)\) into \(2^k\) equally
sized half-open intervals, and define \(f\) as:

\[f_k\pr{x} = \bc
\frac{m}{2^k}&\if 0\le f\pr{x}\le k\text{ and }  f\pr{x}\in \left[\frac{m}{2^k},\frac{m+1}{2^k} \right)\\
\frac{m+1}{2^k}&\if -k\le f\pr{x}\le 0\text{ and }  f\pr{x}\in \left[\frac{m}{2^k},\frac{m+1}{2^k} \right)\\
k&\if f\pr{x}>k\\
-k&\if f\pr{x}<-k
\ec\]

The first two conditions can easily be checked by construction of
\(f_k\)\textquotesingle s. For any \(f\pr{x}\in\br{-k,k}\), we have
\(\abs{f_k\pr x -f\pr x}\le \frac1{2^k}\), for every \(x\in X\) such
that \(f\pr x\in\br{-k,k}\) and for every \(k\). This implies pointwise
convergence that \(\lim_{k\to\infty}f_k\pr x = f\pr x\). Furthermore, if
\(f\) is bounded, then we can remove the condition in the previous
statement that \(f\pr x\in \br{-k,k}\), which means that
\(\abs{f_k\pr x-f\pr x}\le \frac1{2^k}\) for all \(x\in X\) and all
\(k\). Therefore, \(f_1,f_2,\cdots\) converges \emph{uniformly} on \(X\)
to \(f\) if \(f\) is bounded.

\textbf{Remarks:} Be careful about the interpretation of the conclusion
of Luzin\textquotesingle s theorem that \(f|_B\) is a continuous
function on \(B\). This is not the same as saying that \(f\) (on its
original domain) is continuous at each point of \(B\). For example,
\(\chi_{\Q}\) is discontinuous at every point of \(\R\). However,
\(\chi_{\Q}|_{\R\c \Q}\) is a continuous function on \(\R\c \Q\)
(because this function is identically \(0\) on its domain).

\emph{\textbf{Luzin\textquotesingle s Theorem}} Suppose \(g:\R\to \R\)
is a Borel measurable function. Then for every \(\ve>0\), there exists a
closed set \(F\se \R\) such that \(\lm{\R\c F}<\ve\) and \(g|_{F}\) is a
continuous function on \(F\) (\(g|_F:F\to \R\)).

\emph{\textbf{Proof}} The proof goes with the following sketch:

\begin{itemize}
\item
  Prove Luzin\textquotesingle s theorem for simple function \(g_k\), and
  \(g_k\to g\). (pointwise convergence)
\item
  Use Egorov\textquotesingle s theorem to make the previous pointwise
  convergence to uniform convergence.
\item
  Uniform limit of continuous functions is continuous.
\end{itemize}

First prove Luzin\textquotesingle s theorem for simple functions. Assume
\(g\) to be a simple function that

\[g = d_1\chi_{D_1}+\cdots d_n\chi_{D_n}\]

for some \(d_1,d_2,\cdots,d_n\ne 0\) and some disjoint
\(D_1,D_2,\cdots,D_n\se \R\).

Let \(\ve>0\). For each \(k\), there exists a closed set \(F_k\se D_k\),
and an open set \(G_k\supseteq D_k\) such that

\[\lm{D_k\c F_k}\le\frac{\ve}{2n}\\
\lm{G_k\c D_k}\le\frac{\ve}{2n}\]

Because \(G_k\c F_k = \pr{G_k\c D_k}\cup \pr{D_k\c F_k}\), by
subadditivity we have

\[\lm{G_k\c F_k}\le \frac{\ve}{2n}+ \frac{\ve}{2n}= \frac{\ve}{n},\forall k\in \brace{1,2,\cdots,n}\]

Note that simple function must be continuous. Let

\[F = \pr{\bigcup_{k = 1}^nF_k}\cup \pr{\bigcap_{k = 1}^n\pr{\R\c G_k}}\]

Because \(F_k\se D_k\), \(g\) is identically \(d_k\) on \(F_k\). Then
\(g|_{F_k}\) is constant and then continuous for each
\(k\in \brace{1,2,\cdots,n}\), and \(g|_{\R\c G_k}\) is \(0\) and then
continuous for each \(k\). In conclusion, \(g|_F\) is continuous.

By theorem of approximation by simple functions, there exists a sequence
of simple functions \(g_1,g_2,\cdots:\R\to \R\) such that the sequence
converges pointwise to \(g\). Let \(\ve>0\). By the first case above,
for each \(k>0\), there exists a closed set \(C_k\se \R\) such that

\begin{itemize}
\item
  \(\lm{\R\c C_k}\le \frac{\ve}{2^{k+1}}\).
\item
  \(g_k|_{C_k}\) is continuous.
\end{itemize}

Let \(C = \iinf C_k\). Then \(C\) is a closed set and \(C\se C_k\). From
this we can see that \(g_k|_C\) is continuous for every \(k\). By De
Morgan\textquotesingle s law:

\[\lm{\R\c C}\le \sinf\lm{\R\c C_k}\le \sinf \frac{\ve}{2^{k+1}} = \ve\]

For each \(m\in \Z\), consider \(\pr{m,m+1}\). The sequence
\(g_1|_{\pr{m,m+1}},g_2|_{\pr{m,m+1}},\cdots\) converges pointwise to
\(g|_{\pr{m,m+1}}\). This is obviously true because
\(g_k\)\textquotesingle s converge pointwise to \(g\) on \(\R\), so is
the case when restricted to \(\pr{m,m+1}\). By Egorov\textquotesingle s
theorem, there exists a Borel set \(E_m\) such that

\[\lm{\pr{m,m+1}\c E_n}\le \frac{\ve}{2^{\abs{m}+3}}\]

and such that the sequence \(g_1,g_2,\cdots\) converges uniformly to
\(g\) on \(E_m\). Then we have \(g_1,g_2,\cdots\) converges uniformly to
\(g\) on \(C\cap E_m\). Because each \(g_k\) is continuous,
\(g|_{C\cap E_m}\) is continuous. Define

\[D = \bigcup_{m\in \Z}\pr{C\cap E_m}\]

Then \(g|_D\) is continuous. The last step is to show that

\[\lm{\R\c D}\le \cdots \le \ve\]

There exists a closed set \(F\se D\) such that

\[\lm{D\c F}\le \ve-\lm{\R\c D}\]

Then

\[\lm{\R\c F}\le \lm{D\c F}+\lm{\R\c D}\le \ve\]

\(g|_F\) is continuous because \(g|_D\) is continuous.

\textbf{Remarks}

\begin{itemize}
\item
  Every time you come across the case of \(A\se B\), you may not write
  \(\lm{B\c A} = \lm{A}-\lm{B}\).
\item
  Any intersection of closed set, regardless of finitely or infinitely
  many, must be closed. However this is not true for union operation.
\end{itemize}

\begin{center}\rule{0.5\linewidth}{0.5pt}\end{center}

Luzin\textquotesingle s theorem also has a second version.

\emph{\textbf{Theorem (Continuous Extensions of Continuous Functions)}}
Every continuous function on a closed subset of \(\R\) can be extended
to a continuous function on all of \(\R\). More precisely, if
\(F\se \R\) is closed and \(g:F\to \R\) is continuous, then there exists
a continuous function \(h:\R\to \R\) such that \(h|_F = g\).

\emph{\textbf{Luzin\textquotesingle s Theorem (Second Version)}} Suppose
\(E\se \R\) and \(g:E\to \R\) is a Borel measurable function. Then for
every \(\ve>0\), there exists a closed set \(F\se E\) and a continuous
function \(h:\R\to \R\) such that \(\om{E\c F}<\ve\) and
\(h|_F = g|_F\).

\emph{\textbf{Proof}} Suppose \(\ve>0\). Extend \(g\) to a function
\(\tilde g:\R\to \R\) by defining

\[\tilde g\pr{x} = \bc
g\pr{x},\if x\in E\\
0,\if x\in \R\c E
\ec\]

By the first version of Luzin\textquotesingle s theorem, there is a
closed set \(C\se \R\) such that \(\om{\R\c C}<\ve\) and \(\tilde g|_C\)
is a continuous function on \(C\). There exists a closed set
\(F\se \pr{C\cap E}\) such that

\[\om{\pr{\C\cap E}\c F}<\ve-\om{\R\c C}\]

Thus,

\[\om{E\c F}\le \om{\pr{\pr{C\cap E}\c F}\cup \pr{\R\c C}}\le \om{\pr{C\c E}\c F}+\om{\R\c C}<\ve\]

Now \(\tilde g|_F\) is a continuous function on \(F\). Also,
\(\tilde g|_F = g|_F\) (because \(F\se E\)). Use the theorem of
continuous extensions of continuous functions to extend \(\tilde g|_F\)
to a continuous function \(h:\R\to \R\).

\end{document}
